% concatenated from main.tex
\documentclass[11pt]{article}

\usepackage{amsmath, amssymb}
\usepackage{authblk}
\usepackage{geometry}
\usepackage{hyperref}
\usepackage{graphicx}
\usepackage{epstopdf}
\usepackage{mathrsfs}
\usepackage{amsfonts}
\usepackage{amscd,amsthm,bm,psfrag}
\usepackage[numbers,sort&compress]{natbib}

\geometry{margin=1in}

\setlength{\baselineskip}{16.5pt}
\setlength{\evensidemargin}{-4cm} \setlength{\oddsidemargin}{1mm}
\setlength{\textwidth}{16cm} \setlength{\textheight}{22cm}
\setlength{\headsep}{1.4mm}

\makeatletter
\renewcommand{\@seccntformat}[1]{{\csname the#1\endcsname}{\normalsize.}\hspace{.5em}}

\makeatother
\renewcommand{\thesection}{\normalsize \arabic{section}}
\renewcommand{\refname}{\normalsize \bf{References}}
\renewcommand{\thefootnote}{\fnsymbol{footnote}}


\numberwithin{equation}{section}

\def\theequation{\thesection.\arabic{equation}}
% \def \[{\begin{equation}}
% \def \]{\end{equation}}


\newtheorem{thm}{Theorem}[section]
\newtheorem{defi}[thm]{Definition}
\newtheorem{lem}[thm]{Lemma}
\newtheorem{cor}[thm]{Corollary}
\newtheorem{prop}[thm]{Proposition}
\newtheorem{rem}[thm]{Remark}
\newtheorem{con}[thm]{Conjecture}
\newtheorem{prob}[thm]{Problem}
\newtheorem{cla}{Claim}
\newtheorem{exa}{Example}
%\newtheorem{fac}{Factor}[section]
\newtheorem*{thm*}{Theorem}
\newtheorem*{lem*}{Lemma}
\newtheorem*{prop*}{Proposition}

\newcounter{cases}
\newcounter{subcases}[cases]
\newenvironment{mycase}
{
    \setcounter{cases}{0}
    \setcounter{subcases}{0}
    \newcommand{\case}
    {
        \par\indent\stepcounter{cases}\textbf{Case \thecases.}
    }
    \newcommand{\subcase}
    {
        \par\indent\stepcounter{subcases}\textbf{Subcase \thesubcases:}
    }
}
{
    \par
}
\renewcommand*\thecases{\arabic{cases}}
\renewcommand*\thesubcases{\roman{subcases}}

\newtheorem{problem}[thm]{Problem}
%\newtheorem{defin}[thm]{Definition}
%\newtheorem{defn}[thm]{Definition}
%\newtheorem{definition}[thm]{Definition}
%\newtheorem{notation}[thm]{Notation}
\newtheorem{fact}[thm]{Fact}
%\newtheorem{cor}[thm]{Corollary}
%\newtheorem{corollary}[thm]{Corollary}
\newtheorem{question}[thm]{Question}
%\newtheorem{kov}[thm]{Corollary}
%\newtheorem{example}[thm]{Example}
\newtheorem{que}[thm]{Question}
%\newtheorem{rem}[thm]{Remark}
\newtheorem{clm}[thm]{Claim}
\newtheorem{observation}[thm]{Observation}



\DeclareMathOperator{\Sym}{Sym}
\DeclareMathOperator{\Sep}{Sep}
\DeclareMathOperator{\Age}{Age}
\DeclareMathOperator{\Stab}{Stab}
\DeclareMathOperator{\id}{id}
\DeclareMathOperator{\Inv}{Inv}
\DeclareMathOperator{\Aut}{Aut}
\DeclareMathOperator{\End}{End}
\DeclareMathOperator{\Pol}{Pol}
\DeclareMathOperator{\Proj}{Proj}
\DeclareMathOperator{\Cl}{Cl}
\DeclareMathOperator{\Cll}{Cl_{loc}}
\DeclareMathOperator{\pol}{Pol}
\DeclareMathOperator{\Typ}{Typ}
\DeclareMathOperator{\typ}{Typ}
\DeclareMathOperator{\Type}{Typ}
\DeclareMathOperator{\ran}{ran}
\DeclareMathOperator{\dom}{dom}
\DeclareMathOperator{\cyc}{cycl}
\DeclareMathOperator{\sw}{sw}
\DeclareMathOperator{\sep}{sep}
\DeclareMathOperator{\betw}{betw}
\DeclareMathOperator{\sync}{sync}
\DeclareMathOperator{\link}{link}
\DeclareMathOperator{\wlink}{wlink}
\DeclareMathOperator{\pp}{pp}

\newcommand\ex{\ensuremath{\mathrm{ex}}}
\newcommand\rb{\ensuremath{\mathrm{rb}}}
\newcommand\cA{{\mathcal A}}
\newcommand\cB{{\mathcal B}}
\newcommand\cC{{\mathcal C}}
\newcommand\cD{{\mathcal D}}
\newcommand\cE{{\mathcal E}}
\newcommand\cF{{\mathcal F}}
\newcommand\cG{{\mathcal G}}
\newcommand\cH{{\mathcal H}}
\newcommand\cI{{\mathcal I}}
\newcommand\cJ{{\mathcal J}}
\newcommand\cK{{\mathcal K}}
\newcommand\cL{{\mathcal L}}
\newcommand\cM{{\mathcal M}}
\newcommand\cN{{\mathcal N}}
\newcommand\cP{{\mathcal P}}
\newcommand\cY{{\mathcal Y}}
\newcommand\cQ{{\mathcal Q}}
\newcommand\cR{{\mathcal R}}
\newcommand\cS{{\mathcal S}}
\newcommand\cT{{\mathcal T}}
\newcommand\cU{{\mathcal U}}
\newcommand\cV{{\mathcal V}}
\newcommand\PP{{\mathbb P}}
\renewcommand{\P}{\mathbb P}
\renewcommand{\S}{\mathcal S}

\newcommand{\cl}[1]{\langle #1 \rangle}
\newcommand{\To}{\rightarrow}
\newcommand{\nin}{\notin}
\newcommand{\inv}{^{-1}}
\newcommand{\mult}{\times}

\DeclareMathOperator{\tp}{tp}

\DeclareMathOperator{\Rev}{Rev}
\DeclareMathOperator{\Frlim}{Frlim}
\DeclareMathOperator{\R}{Rev}
\DeclareMathOperator{\T}{Turn}
\DeclareMathOperator{\M}{Max}
\DeclareMathOperator{\pari}{Par}

\def\lc{\left\lceil}
\def\rc{\right\rceil}

\def\lf{\left\lfloor}
\def\rf{\right\rfloor}


%\newcommand{\ex}{\rm ex}
\newcommand{\Berge}{\rm Berge}
%\def\Mod{{\rm Mod}}
%\def \CL{\rm CL}
%\def\TN{\rm TN}
%\def\O{\mathcal{O}(M)}
%\def\AO{A\mathcal{O}(M)}
%\def\Z{{\Bbb Z}}

\newenvironment{kst}
{\setlength{\leftmargini}{1.3\parindent}
 \begin{itemize}
 \setlength{\itemsep}{-1.1mm}}
{\end{itemize}}

\newenvironment{lst}
{\setlength{\leftmargini}{0.75\parindent}
 \begin{itemize}
 \setlength{\itemsep}{-1.1mm}}
{\end{itemize}}

\newenvironment{wst}
{\setlength{\leftmargini}{1.5\parindent}
 \begin{itemize}
 \setlength{\itemsep}{-1.1mm}}
{\end{itemize}}


\renewcommand\Authands{, }


%On Tur\'{a}n problems for families of Berge hypergraphs
\begin{document}
\title{The Tur\'{a}n number of Berge matchings}
% \author{
% Xiamiao Zhao\thanks{\small Al. Email: \small \texttt{g}.}\,, \hspace{0.1em}
% Zixuan Yang\thanks{\small Al. Email: \small \texttt{g}.}\,, \hspace{0.1em}
% Yichen Wang\thanks{\small Al. Email: \small \texttt{g}.}\,, \hspace{0.1em}
% Yuhang Bai\thanks{\small Al. Email: \small \texttt{g}.}\,, \hspace{0.1em}
% Junpeng Zhou\thanks{\small \textit{Corresponding author.} Department of Mathematics, Shanghai University, Shanghai 200444, P.R. China. Email: \small \texttt{junpengzhou@shu.edu.cn}.} \thanks{\small Newtouch Center for Mathematics of Shanghai University, Shanghai 200444, P.R. China.}
% }

\author[1]{\small\bf Yichen Wang\thanks{E-mail: wangyich22@mails.tsinghua.edu.cn}}
\author[2,3]{\small\bf Zixuan Yang\thanks{E-mail: yangzixuan@nwpu.edu.cn}}
\author[1]{\small\bf Xiamiao Zhao\thanks{E-mail: zxm23@mails.tsinghua.edu.cn}}
\author[2,3]{\small\bf Yuhang Bai\thanks{E-mail: yhbai@mail.nwpu.edu.cn}}
\author[4,5]{\small\bf Junpeng Zhou\thanks{\textit{Corresponding author}. E-mail: junpengzhou@shu.edu.cn}}

\affil[1]{\small Department of Mathematical Sciences, Tsinghua University, Beijing 100084, P.R. China.}
\affil[2]{\small School of Mathematics and Statistics, Northwestern Polytechnical University, Xi'an 710129, Shaanxi, P.R. China.}
\affil[3]{\small Xi'an-Budapest Joint Research Center for Combinatorics, Xi'an 710129, Shaanxi, P.R. China.}
\affil[4]{\small Department of Mathematics, Shanghai University, Shanghai 200444, P.R. China.}
\affil[5]{\small Newtouch Center for Mathematics of Shanghai University, Shanghai 200444, P.R. China.}



\date{}

\maketitle

\begin{abstract}
    Given a graph $F$, an $r$-uniform hypergraph $\mathcal{H}$ is a {\em Berge-$F$} if there is a bijection $\phi:E(F)\to E(\mathcal{H})$ such that $e\subseteq \phi(e)$ for each $e\in E(F)$. Given a family $\mathcal{F}$ of $r$-uniform hypergraphs, an $r$-uniform hypergraph is $\mathcal{F}$-free if it does not contain any member of $\mathcal{F}$ as a subhypergraph. The Tur\'{a}n number of $\mathcal{F}$ is the maximum number of hyperedges in an $\mathcal{F}$-free $r$-graph on $n$ vertices. Let $M_{s+1}$ denote a matching of size $s+1$, i.e., the graph consisting of $s+1$ independent edges. Khormali and Palmer [\textit{European J. Combin.} 102 (2022) 103506] completely determined the Tur\'{a}n number of Berge matchings for sufficiently large $n$. Subsequently, Kang, Ni, and Shan [\textit{Discrete Math.} 345 (2022) 112901] determined the exact value of the Tur\'{a}n number of Berge-$M_{s+1}$ for all $n$ when $r \le s-1$ or $r \ge 2s+2$. In this paper, we settle the final open case $s \le r \le 2s+1$, thereby completing the determination of the Tur\'{a}n number of Berge matchings.
    %  Moreover, we establish several exact and general results on the Tur\'{a}n numbers of Berge matchings together with a single $r$-graph, as well as of Berge matchings together with Berge bipartite graphs.
    %  Finally, we generalize the results on Tur\'{a}n problems for Berge hypergraphs proposed by Gerbner, Methuku, and Palmer [\textit{Eur. J. Comb. 86 (2020) 103082}]. %are generalized.
\end{abstract}

{\noindent{\bf Keywords}: Tur\'{a}n number, Berge hypergraph, matching}

{\noindent{\bf AMS subject classifications:} 05C35, 05C65}


\section{\normalsize Introduction}

% A \textit{hypergraph} $\cH=(V(\cH),E(\cH))$ consists of a vertex set $V(\cH)$ and a hyperedge set $E(\cH)$, where each hyperedge in $E(\cH)$ is a nonempty subset of $V(\cH)$. If $|e|=r$ for every $e\in E(\cH)$, then $\cH$ is called an \textit{$r$-uniform hypergraph} ($r$-graph for short). For simplicity, let $e(\cH):=|E(\cH)|$. The \textit{degree} $d_\cH(v)$ of a vertex $v$ is the number of hyperedges containing $v$ in $\cH$.

%In this paper, we use capital letter to represent graphs~(e.g., $F$, $G$), calligraphic capital letter to represent hypergraphs~(e.g., $\cG$, $\cH$).
%  frank letter to represent class of hypergraphs~(e.g., $\mathfrak{F}$). For a graph $F$, we use $V(F)$ to denote its vertex set, $E(F)$ to denote its edge set, and $e(F)$ to denote the number of edges, i.e., $e(F) = |E(F)|$. Similarly, for a hypergraph $\cH$, we define $V(\cH)$, $E(\cH)$ and $e(\cH)$ in the same way.

A \textit{hypergraph} $\cH$, denoted by $\cH=(V(\cH),E(\cH))$, on a finite vertex set $V(\cH)$ is a family $E(\mathcal{H})$ of subsets of $V(\cH)$, called \textit{hyperedges}.
For an integer $r\geq2$, a hypergraph $\cH$ is called an \textit{$r$-uniform hypergraph} ($r$-graph for short) if every hyperedge of $\cH$ contains exactly $r$ vertices. %We identify a hypergraph $\cH$ with its hyperedge set $E(\cH)$ and denote its vertex set by $V(\cH)$.
The size of $E(\cH)$ is denoted by $e(\cH)$. The \textit{degree} of a vertex $v$ in $\cH$, denoted $d_{\cH}(v)$, is the number of hyperedges of $\cH$ that contain $v$.
When there is no ambiguity regarding $\cH$, we simply write $d(v)$ instead of $d_\cH(v)$.
Throughout this paper, %we use capital letter to represent graphs (e.g., $F$, $G$), and
we assume that all $r$-graphs are simple, i.e., no loops and no multiple edges (hyperedges).

Given a family $\cF$ of $r$-graphs, we say $\cH$ is \textit{$\cF$-free} if $\cH$ does not contain any member of $\cF$ as a subhypergraph. The \textit{Tur\'{a}n number} ${\rm{ex}}_r(n,\cF)$ of $\cF$ is the maximum number of hyperedges in an $\cF$-free $r$-graph on $n$ vertices. When $r=2$, we write ${\rm{ex}}(n,\cF)$ instead of ${\rm{ex}}_2(n,\cF)$.
Tur\'{a}n problems on graphs and hypergraphs are central topics in extremal combinatorics. A classical result in extremal graph theory is the Erd\H{o}s-Gallai theorem \cite{gallai1959maximal}, which determines the exact Tur\'{a}n number of matchings. 

It is natural to investigate the Tur\'{a}n problem for matchings in hypergraphs. For a graph $F$ and an $r$-graph $\cF$, we say $\cF$ is a \textit{Berge copy} of $F$~(a Berge-$F$ for short) if $V(F) \subseteq V(\cF)$ and there is a bijection $\phi: E(F) \rightarrow E(\cF)$ such that $e \subseteq \phi(e)$ for each $e \in E(F)$. The graph $F$ is called a \textit{core} of $\cF$. Observe that for a fixed graph $F$ there are many hypergraphs that are Berge copies of $F$. For convenience, we refer to this collection of hypergraphs as $\mathcal{B}F$. 
In 1989, Berge~\cite{berge1984hypergraphs} introduced the notion of a \textit{Berge cycle}. Gy\H{o}ri, Katona and Lemons~\cite{GYORI2016238} defined the notion of \textit{Berge paths} and generalized the Erd\H{o}s-Gallai theorem to Berge paths in 2016. Later, Gerbner and Palmer \cite{gerbner2017extremal} generalized the established notions of Berge cycle and Berge path to general graphs.
Tur\'{a}n problems on Berge hypergraphs have been extensively studied, yielding numerous related results.
For a short survey on Tur\'{a}n problems on Berge hypergraphs, one can refer to Subsection 5.2.2 in~\cite{gerbner2018extremal}. 
In particular, Khormali and Palmer~\cite{KHORMALI2022103506} completely determined the Tur\'{a}n number of Berge matchings for sufficiently large $n$.
Let $M_{s+1}$ denote a matching of size $s+1$, i.e., the graph consisting of $s+1$ independent edges. 

\begin{thm}[Khormali and Palmer~\cite{KHORMALI2022103506}]\label{thm: Khormali Palmer}
Fix integers $s \geq 1$ and $r \geq 2$. Then for $n$ large enough,
\begin{eqnarray*}
{\rm{ex}}_r(n,\mathcal{B}M_{s+1})=
\begin{cases}
s , & {\rm{if}}\ r \geq 2s + 1; \\
\binom{2s+1}{r}, & {\rm{if}}\ s+1 < r < 2s + 1; \\
n - s , & {\rm{if}}\ r = s+1; \\
\binom{s}{r-1} (n - s ) + \binom{s}{r}, & {\rm{if}}\ r \leq s.
\end{cases}
\end{eqnarray*}
\end{thm}

Note that if $n \le 2s+1$, it is trivial that ${\rm{ex}}_r(n,\mathcal{B}M_{s+1}) = \binom{n}{r}$. For $n\ge 2s+2$ and $r\geq3$, Kang, Ni and Shan \cite{KANG2022112901} determined the exact value of $\ex_r(n,\mathcal{B}M_{s+1})$ in the case when $r\leq s-1$ or $r\geq 2s+2$. However, the case $s\leq r\leq 2s+1$ remains open.

\begin{thm}[Kang, Ni and Shan~\cite{KANG2022112901}]\label{thm: Kang Ni Shan}
    Fix integers $s \ge 1$ and $r \ge 3$. For any $n \ge 2s+2$.
    \begin{eqnarray*}
    {\rm{ex}}_r(n,\mathcal{B}M_{s+1})=
    \begin{cases}
    \max\{ \binom{2s+1}{r}, \binom{s}{r-1} (n - s ) + \binom{s}{r}\}, & {\rm{if}}\ r \leq s - 1; \\
    s & {\rm{if}}\ r \ge 2s+2.
    \end{cases}
    \end{eqnarray*}
\end{thm}

The purpose of this paper is to solve the remaining case of the Tur\'{a}n number of Berge matchings as stated in Theorem~\ref{thm: exact berge}. 
For non-negative integers $a$ and $b$, let $\binom{a}{b}$ denote the binomial coefficient. In particular, we define $\binom{a}{b}=0$ if $a<b$. %Here for positive integers $a$ and $b$, define $\binom{a}{b}=0$ if $a<b$.
% when $s\leq r\leq 2s+1$.

\begin{thm}\label{thm: exact berge}
    Let integers $s \ge 1$ and $r \ge 2$. For all $n \ge 2s+2$,
    \begin{eqnarray*}
    {\rm{ex}}_r(n,\mathcal{B}M_{s+1})=
    \begin{cases}
    \max\left\{ \binom{2s+1}{r}, \binom{s}{r-1} (n - s ) + \binom{s}{r}\right\}, & {\rm{if}}\ r \leq s + 1; \\
    \binom{2s+1}{r} , & {\rm{if}}\ s + 2 \leq r \leq 2s; \\
    s, & {\rm{if}}\ r \ge 2s+1.
    \end{cases}
    \end{eqnarray*}
\end{thm}

It is worth mentioning that the Tur\'{a}n number and the generalized Tur\'{a}n number of a matching together with another graph have been widely studied~(see~\cite{ALON2024223,Gerbner2024matching, 2023arXiv230105625M},\cite{Zhu}).

The $r$-expansion of a graph $F$ is a special member of $\mathcal{B}F$, defined by enlarging each edge of $F$ with a set of $r-2$ new vertices such that all these new sets are disjoint.
When considering the Tur\'{a}n number of the expansion of a matching, one obtains the famous Erd\H{o}s Matching Conjecture, which has been solved for sufficiently large $n$~\cite{FRANKL20131068}.
% When considering the Tur\'{a}n number of the expansion of a matching, it is the famous Erd艖s' Matching Conjecture, which is solved when $n$ is sufficiently large~\cite{FRANKL20131068}.
There are also many result on the Tur\'{a}n number of the expansion of a matching together with the expansion of another graph.
The reader may refer to~\cite{GERBNER2025104155, 2025arXiv250704579W, 2025arXiv251121096Y, 2025arXiv251117000C, zhou2026linear}.
The above results reveal the importance of matchings in extremal combinatorics.


% \section{\normalsize Main results}\label{sec: main results}
% %\subsection{\normalsize Tur\'{a}n number of Berge matchings}


% Our first result is to determine the exact value of ${\rm{ex}}_r(n,\mathcal{B}M_{s+1})$ for all $n$ in the case $s\leq r\leq 2s+1$ in Theorem~\ref{thm: exact berge}, which will be proved in Section~\ref{sec: berge exact}.
% For positive integers $a$ and $b$, define $\binom{a}{b}=0$ if $a<b$. %when $s\leq r\leq 2s+1$.

% \begin{thm}\label{thm: exact berge}
%     Let integers $s \ge 1$ and $r \ge 2$. For all $n \ge 2s+2$,
%     \begin{eqnarray*}
%     {\rm{ex}}_r(n,\mathcal{B}M_{s+1})=
%     \begin{cases}
%     \max\left\{ \binom{2s+1}{r}, \binom{s}{r-1} (n - s ) + \binom{s}{r}\right\}, & {\rm{if}}\ r \leq s + 1; \\
%     \binom{2s+1}{r} , & {\rm{if}}\ s + 2 \leq r \leq 2s; \\
%     s, & {\rm{if}}\ r \ge 2s+1.
%     \end{cases}
%     \end{eqnarray*}
% \end{thm}



% Moreover, we provide a formula for the Tur\'{a}n number of Berge matching together with a single $r$-graph $\cF$, namely $\ex_r(n,\mathcal{B}M_{s+1}\cup\{\cF\})$, and determine the asymptotics of the Tur\'{a}n number of Berge matchings together with Berge-$G$ where $G$ is a bipartite graph, namely $\ex_r(n,\mathcal{B}M_{s+1}\cup\mathcal{B}G)$.
% Finally, we extend two general lemmas for Tur\'{a}n problems on Berge hypergraphs from~\cite{GERBNER2020103082}, and present some applications.
% %Also, we give some exact results similar to Theorem~\ref{thm: gerbner matching} for Tur\'{a}n numbers of a Berge matching and one hypergraph, i.e., ${\rm{ex}}_r(n,\mathcal{B}M_{s+1},\{\cF\})$.
% %Moreover, we consider the Tur\'{a}n number of a Berge matching and a Berge bipartite graph, i.e., ${\rm{ex}}_r(n,\mathcal{B}M_{s+1},\mathcal{B}G)$ where $G$ is a bipartite graph.



% \subsection{\normalsize Berge matchings and a single $r$-graph}

% We first introduce some definitions from~\cite{GERBNER2025104155}.
% A \textit{proper $k$-coloring} of $\cH$ is a mapping from $V(\cH)$ to a set of $k$ colors such that no hyperedge is monochromatic. The \textit{chromatic number} $\chi(\cH)$ of $\cH$ is the minimum colors needed for proper coloring $\cH$.
% % Note that the chromatic number of the Fano plane is 3, and the chromatic number of the $r$-expansion of any graph is 2 for all integers $r \geq 3$.

% If $\chi(\cH)=2$, then we take the colors of a proper 2-coloring of $\cH$ to be red and blue.
% In the following, when we say \textit{red-blue coloring}, we always mean a proper $2$-coloring. Let $p(\cH)$ denote the minimum number of red vertices in a red-blue coloring of $\cH$. We say that a red-blue coloring of $\cH$ is \textit{strong} if each hyperedge contains exactly one red vertex. We remark that the set of red vertices in a strong red-blue coloring is also called a \textit{crosscut} in the literature.
% Let $q(\cH)$ denote the minimum number of red vertices in a strong red-blue coloring of $\cH$, with $q(\cH)=\infty$ if $\cH$ does not have a strong red-blue coloring. Clearly, $p(\cH) \leq q(\cH)$.


% A \textit{strongly independent set} of a hypergraph $\cH$ is a subset of $V(\cH)$ that intersects every hyperedge in at most one vertex. Note that in a strong red-blue coloring, the set of red vertices is a strongly independent set. Our main result in this subsection is as follows.

% %Observe that a nonempty $r$-graph $\cH$ contains a strongly independent set such that deleting it results in the empty graph if and only if $\chi(\cH)=2$ and $q(\cH)<\infty$. Therefore, a nonempty $r$-graph $\cH$ has the property that deleting any strongly independent set leaves it nonempty if and only if $\chi(\cH)>2$, or $\chi(\cH)=2$ and $q(\cH)=\infty$.


% \begin{thm}\label{thm: berge matching and a hypergraph}
% Let integers $2\leq r\leq s+1$ and $n$ be sufficiently large. Suppose $\cF$ is an $r$-graph with $\chi(\cF)>2$, or $\chi(\cF)=2$ and $q(\cF)=\infty$. %$F$ is a nonempty $r$-graph such that deleting any strongly independent set leaves it nonempty.
% Then
% \[
%     \ex_r(n,\mathcal{B}M_{s+1}\cup\{\cF\})=\ex_r(s,\mathcal{D}(\cF))+ \binom{s}{r-1}(n-s),
% \]
% where $\mathcal{D}(\cF)$ is the family of $r$-graphs obtained by deleting a strongly independent set from $\cF$.
% \end{thm}


% \begin{rem}
%     Note that in the above theorem we assume that $r\leq s+1$.
%     If $r\geq s+2$, then $\ex_r(n,\mathcal{B}M_{s+1}\cup\{\cF\})\leq \ex_r(n,\mathcal{B}M_{s+1})\leq \binom{2s+1}{r}$ by Theorem~\ref{thm: exact berge}.
%     This implies that $\ex_r(n,\mathcal{B}M_{s+1}\cup\{\cF\})$ is a constant which does not depend on $n$ when $r\ge s+2$.
% \end{rem}

% Theorem~\ref{thm: berge matching and a hypergraph} can be viewed as a Berge-version extension of Theorem~\ref{thm: gerbner matching}.
% We first provide two applications of Theorem~\ref{thm: berge matching and a hypergraph}. For $k>r$, let $\cK_k^{(r)}$ denote the complete $r$-graph on $k$ vertices. We color the vertices of $\cK_k^{(r)}$. It is easy to see that the number of vertices of each color is at most $r-1$. Otherwise, a monochromatic hyperedge would exist. Then $\chi(\cK_k^{(r)})=\lceil\frac{k}{r-1}\rceil$.
% When $k > 2(r-1)$, we have $\chi(\cK_k^{(r)}) \geq 3$.
% When $k\leq 2(r-1)$, we claim that $q(\cK_k^{(r)})=\infty$.
% Indeed, if $q(\cK_k^{(r)})<\infty$, then $q(\cK_k^{(r)})=1$ as any two vertices in $\cK_k^{(r)}$ are contained in some hyperedge. However, after deleting one red vertex, there are still at least $r$ blue vertices, which implies the existence of a monochromatic hyperedge, a contradiction. Thus, by Theorem~\ref{thm: berge matching and a hypergraph}, we derive the following corollary.

% \begin{cor}
% Let integers $2\leq r\leq s+1$, $k\geq r+1$ and $n$ be sufficiently large. Then
% \[
%     \ex_r(n,\mathcal{B}M_{s+1}\cup\{\cK_k^{(r)}\})=\ex_r(s,\cK_{k-1}^{(r)})+\binom{s}{r-1}(n-s).
% \]
% \end{cor}

% Observe that if $k>s+1$, then $\ex_r(s,\cK_{k-1}^{(r)})=\binom{s}{r}$ and $\ex_r(n,\mathcal{B}M_{s+1}\cup\{\cK_k^{(r)}\})=\binom{s}{r}+\binom{s}{r-1}(n-s)$.
% % If $k\geq 2(s+1)$, then $\ex_r(n,\mathcal{B}M_{s+1}\cup\{\cK_k^{(r)}\})=\ex_r(n,\mathcal{B}M_{s+1})$ as $\cK_k^{(r)}$ contains a $\mathcal{B}M_{s+1}$.


% The other application is about the Fano plane.
% Let $\cF_7$ denote the Fano plane. Since any two vertices in $\cF_7$ are contained in some hyperedge, any single vertex in $\cF_7$ forms a maximum strongly independent set.
% By deleting any vertex from $\cF_7$, we obtain the $3$-graph $\mathcal{F}_6$ with vertex set $\{1,2,\dots,6\}$ and hyperedge set $\big\{123,145,246,356\big\}$.
% See Figure~\ref{fig: fano} for an illustration of $\cF_7$ and $\cF_6$.
% Note that $\chi(\cF_7)=3$. Therefore, Theorem~\ref{thm: berge matching and a hypergraph} implies the following result.

% \begin{figure}[t]
%     \centering
%     \includegraphics[width=0.5\textwidth]{./src/Berge-02.jpg}%
%     \caption{The Fano plane $\cF_7$ and the hypergraph $\cF_6$.}\label{fig: fano}
% \end{figure}

% \begin{cor}
% Let integers $s\geq 2$ and $n$ be sufficiently large. Then
% \[
%     \ex_3(n,\mathcal{B}M_{s+1}\cup\{\cF_7\})=\ex_3(s,\cF_6)+\binom{s}{2}(n-s).
% \]
% \end{cor}

% The proof of Theorem~\ref{thm: berge matching and a hypergraph} will be presented in Section~\ref{sec: berge matching and a hypergraph}.

% \subsection{\normalsize Berge matchings and Berge bipartite graphs}\label{subsec: berge matching and berge bigraph}

% Now let us consider forbidding Berge matchings together with Berge-$G$ where $G$ is a bipartite graph.
% We have the following theorem.
% % Let $A$ and $B$ be the bipartition $A$ and $B$ of $G$.
% % Assume $|A| \leq |B|$.
% % and let $\cH$ be an $r$-uniform $\mathcal{B}G$. By coloring $A$ red and $V(\cH)\backslash A$ blue, it can be seen that $\cH$ has a proper 2-coloring. Therefore, $\chi(\cH)=2$ for any $\cH\in \mathcal{B}G$. We obtain the following result on Berge bipartite graphs.

% \begin{thm}\label{thm: berge matching and berge bigraph}
%     Let $s$ and $r\geq2$ be fixed integers. Let $G$ be a bipartite graph with $p=p(G)$ and $n$ sufficiently large.

%     \textbf{(i)} If $r< p\leq s$, then $\ex_r(n,\mathcal{B}M_{s+1}\cup\mathcal{B}G)= \binom{p-1}{r-1}n+O(1)$.

%     %\textbf{(ii)} If $p\leq r\leq s$, then......

%     \textbf{(ii)} If $p\leq s<r-1$, then $\ex_r(n,\mathcal{B}M_{s+1}\cup\mathcal{B}G)\leq\max\big\{s,\binom{2s+1}{r}\big\}$.

%     \textbf{(iii)} If $p>s$, then $\ex_r(n,\mathcal{B}M_{s+1}\cup\mathcal{B}G)= \ex_r(n,\mathcal{B}M_{s+1})$.
% \end{thm}


% In the above theorem, the case $p \leq r \leq s+1$ remains unresolved.
% For the complete bipartite graph $K_{p,q}$ with $p \leq q$, we obtain the asymptotics for $\ex_r(n, \mathcal{B}M_{s+1} \cup \mathcal{B}K_{p,q})$ in that case.
% %In the above theorem, we miss the case when $p\leq r \leq s+1$.
% %When $G$ is a complete bipartite graph $K_{p,q}$~($p \le q$), we give some asymptotics of $\ex_r(n,\mathcal{B}M_{s+1}\cup\mathcal{B}K_{p,q})$ in this hard case.

% % Let $K_{p,q}$ be the complete bipartite graph with parts of size $p$ and $q$, where $p\leq q$.
% % We determine the asymptotics of $\ex_r(n,\mathcal{B}M_{s+1}\cup\mathcal{B}K_{p,q})$ when $p\leq r \leq s$.
% The \textit{$2$-shadow} $\partial\cH$ of a hypergraph $\cH$ is the graph consisting of all $2$-sets contained in some hyperedge of $\cH$.
% For $r\geq3$, let $\cG_s^{r-1}$ be an $s$-vertex $(r-1)$-graph with the maximum number of hyperedges such that $\partial\cG_s^{r-1}$ is $K_{p,q}$-free.
% Note that when $r=p+q\leq s$, we have $e(\cG_s^{r-1})\geq 1$, since the $2$-shadow of a single $(r-1)$-edge is $K_{p,q}$-free.

% \begin{thm}\label{thm: berge matching and berge Kpq}
%     Fix integers $s$ and $r\geq2$. Let $n$ be sufficiently large.

%     \textbf{(i)} If $p\leq r \leq \min\{s,p+q-1\}$, then
%     \[
%         \ex_r(n,\mathcal{B}M_{s+1}\cup\mathcal{B}K_{p,q})=(p-1)n+O(1).
%     \]

%     \textbf{(ii)} If $r=p+q \leq s$, then
%     \[
%     \ex_r(n,\mathcal{B}M_{s+1}\cup\mathcal{B}K_{p,q}) = \min\{p-1,e(\cG_s^{r-1})\}\cdot n+O(1).
%     \]
%     % Moreover, if $(p-1)(p+q-1)\leq s$, then $$\ex_r(n,\mathcal{B}M_{s+1}\cup\mathcal{B}K_{p,q})=(p-1)n+O(1).$$

%     \textbf{(iii)} If $p+q<r\leq s$, then
%     \[
%         \ex_r(n,\mathcal{B}M_{s+1}\cup\mathcal{B}K_{p,q})\leq s+\binom{2s}{r}+\binom{2s}{r-1}(pq-1).
%     \]
% \end{thm}

% For the case $p\leq r=s+1$, we have the following result.

% \begin{prop}\label{prop2.8}
%     Let integers $r\geq2$ and $n$ be sufficiently large.

%     \textbf{(i)} If $p=r=s+1$, then $\ex_r(n,\mathcal{B}M_{s+1}\cup\mathcal{B}K_{p,q})=\ex_r(n,\mathcal{B}M_{s+1})$.

%     \textbf{(ii)} If $p<r=s+1\leq p+q$, then $\ex_r(n,\mathcal{B}M_{s+1}\cup\mathcal{B}K_{p,q})= n+O(1)$.

%     \textbf{(iii)} If $p+q<r=s+1$, then $\ex_r(n,\mathcal{B}M_{s+1}\cup\mathcal{B}K_{p,q})\leq s+\binom{2s}{r}+\binom{2s}{r-1}(pq-1)$.
% \end{prop}

% %We remark that Tur\'{a}n problems on Berge bipartite graphs have been studied previously. In particular, Gerbner, Methuku, and Vizer~\cite{gerbner2019asymptotics} determined the asymptotics of the Tur\'{a}n number of Berge-$K_{2,q}$ in $3$-graphs, and proved general upper and lower bounds for the Tur\'{a}n number of Berge-$K_{2,q}$ in $r$-graphs.

% The proofs of Theorems~\ref{thm: berge matching and berge bigraph} and \ref{thm: berge matching and berge Kpq}, and Proposition~\ref{prop2.8} will be presented in Section~\ref{sec: berge matching and berge bigraph}.

% \subsection{\normalsize Finite families of Berge hypergraphs} %General results for Tur\'{a}n problems on Berge hypergraphs

% In this subsection, we use a slightly different notation.
% In the Section~\ref{subsec: berge matching and berge bigraph}, we consider the Tur\'{a}n problems on two Berge families.
% It is natural to ask what happens for more general families of Berge hypergraphs.
% Let $\mathcal{F}$ be a finite family of graphs and $\mathcal{B} \mathcal{F}:=\bigcup_{F\in \mathcal{F}}\mathcal{B}F$. Now we consider Tur\'{a}n problems on $\mathcal{B}\mathcal{F}$.


% Given a graph $G$, a \textit{red-blue graph} $G_{{\rm red{\text -}blue}}$ is formed by coloring the edges of $G$ with two distinct colors, namely red and blue. For a red-blue graph $G_{{\rm red{\text -}blue}}$, let $G_{{\rm red}}$ denote the subgraph spanned by the red edges of $G_{{\rm red{\text -}blue}}$, and let $G_{{\rm blue}}$ denote the subgraph spanned by the blue edges of $G_{{\rm red{\text -}blue}}$. Set $\varepsilon(r,G_{{\rm red{\text -}blue}}):=e(G_{{\rm red}})+\mathcal{N}(K_r,G_{{\rm blue}})$.


% Inspired by the results of Gerbner, Methuku and Palmer in~\cite{GERBNER2020103082}, we present the following more generalized results.

% \begin{lem}\label{lem3}
% Let $\mathcal{F}$ be a finite family of graphs. Then
% \begin{eqnarray*}
% {\rm{ex}}_r(n,\mathcal{B}\mathcal{F})\leq \max\Big\{\varepsilon(r,G_{{\rm red{\text -}blue}})\,\big|\, G\ {\rm is}\ {\rm an}\ n{\rm {\text -}vertex}\ \mathcal{F}{\rm {\text -}free}\ {\rm graph}\Big\}.
% \end{eqnarray*}
% \end{lem}


% Let
% \begin{eqnarray*}
% \mathcal{F}^-=\Big\{F^-\,\big|\, F^-\ {\rm is\ a\ graph\ obtained\ by\ deleting\ a\ vertex\ from}\ F,\ F\in \mathcal{F}\Big\}.
% \end{eqnarray*}

% By Lemma~\ref{lem3}, we obtain a general upper bound for the Tur\'{a}n number of $\mathcal{B}\mathcal{F}$.

% \begin{thm}\label{lem4}
% Let $\mathcal{F}$ be a finite family of graphs and $f := f(n)$ such that $ {\rm ex}(n,K_{r-1},\mathcal{F}^-) \le f n$ for every $n$, then
% \begin{eqnarray*}
% {\rm{ex}}_r(n,\mathcal{B}\mathcal{F})\leq \max\Big\{\frac{2f}{r},1\Big\}\cdot{\rm{ex}}(n,\mathcal{F}).
% \end{eqnarray*}
% \end{thm}


% Next, we present an application of Theorem~\ref{lem4}.
% Let $\mathcal{F}=\big\{M_{s+1},K_{k+1}\big\}$ and $\mathcal{F}^-=\big\{M_{s}\cup\{v\},K_{k}\big\}$.
% For a matching $M_s$ and a vertex $v\notin V(M_s)$, it is easy to see that if $n\geq 2s+1$, then a $n$-vertex graph is $M_s\cup\{v\}$-free if and only if it is $M_s$-free. Thus, ${\rm ex}(n,K_{r-1},\{M_{s}\cup\{v\},K_{k}\})={\rm ex}(n,K_{r-1},\{M_{s},K_{k}\})$ for $n\geq 2s+1$.
% By Theorems~\ref{thm: Alon frankl}, \ref{lem12} and \ref{lem4}, we obtain the following corollary.

% \begin{cor}\label{cor3}
% Let $n \geq 2s+1$ and $k \geq r \geq 3$. Then
% \begin{eqnarray*}
% {\rm{ex}}_r(n,\mathcal{B}\{M_{s+1},K_{k+1}\})\leq Z(n,r,s,k)\cdot\max\big\{e(T(2s+1,k)),e(G(n,k,s))\big\},
% \end{eqnarray*}
% where $Z(n,r,s,k):=\max\Big\{\frac{2\mathcal{N}(K_{r-1},T(2s-1,k-1))}{rn},\frac{2\mathcal{N}(K_{r-1},G(n,k-1,s-1))\big\}}{rn},1\Big\}$.
% \end{cor}

% Note that if $n\leq 2s+1$, then ${\rm{ex}}_r(n,\mathcal{B}\{M_{s+1},K_{k+1}\})={\rm{ex}}_r(n,\mathcal{B}K_{k+1})$. If $k\geq 2s+1$, then $M_{s+1}\subseteq K_{k+1}$. This implies that ${\rm{ex}}_r(n,\mathcal{B}\{M_{s+1},K_{k+1}\})={\rm{ex}}_r(n,\mathcal{B}M_{s+1})$ for $k\geq 2s+1$.

% The proof of Lemma~\ref{lem3} and Theorem~\ref{lem4} will be presented in Section~\ref{sec: lems}.

% -------------------------------
\section{\normalsize Proof}\label{sec: berge exact}

Given an $r$-graph $\cH$, the \textit{neighborhood} $N_\cH(v)$ of a vertex $v$ in $\cH$ is the set of all vertices adjacent to $v$ in $\cH$. For a set $S\subseteq V(\cH)$, the \textit{neighborhood} of $S$ is defined as $N_\cH(S)=\bigcup_{v \in S} N_\cH(v)\setminus S$.
Recall that $d_{\cH}(v)$ denotes the \textit{degree} of a vertex $v$ in $\cH$.
%When $S$ contains only one vertex $v$, we simply write $N_\cH(v)$.
%And we use $d_\mathcal{H}(v)$ to denote the degree of $v$ in $\mathcal{H}$, i.e., $d_\mathcal{H}(v) = |N_\mathcal{H}(v)|$.
Let $\cH[U]$ denote the subhypergraph of $\cH$ induced by $U\subseteq V(\cH)$.
We say that a Berge matching $\cM$ in $\cH$ is \textit{maximum} if $\cH$ contains no Berge matching with a larger size than $\cM$.
When we say that $\cH$ contains a $\mathcal{B}M_{s+1}$, we mean that $\cH$ contains a Berge copy of $M_{s+1}$.
% When referring to $\cH$ contains a $\cB M_{s+1}$, we mean $\cH$ contains a Berge copy of $M_{s+1}$. %one of the hypergraphs in the family $\cB M_{s+1}$.

% Now let us start with the proof of Theorem~\ref{thm: exact berge} that we restate here for convenience.


% \begin{thm*}
%     Let integers $s \ge 1$ and $r \ge 2$. For all $n \ge 2s+2$,
%     \begin{eqnarray*}
%     {\rm{ex}}_r(n,\mathcal{B}M_{s+1})=
%     \begin{cases}
%     \max\left\{ \binom{2s+1}{r}, \binom{s}{r-1} (n - s ) + \binom{s}{r}\right\}, & {\rm{if}}\ r \leq s + 1; \\
%     \binom{2s+1}{r} , & {\rm{if}}\ s + 2 \leq r \leq 2s; \\
%     s, & {\rm{if}}\ r \ge 2s+1.
%     \end{cases}
%     \end{eqnarray*}
% \end{thm*}

Note that for the cases $r \le s-1$ and $r \ge 2s+2$ in Theorem \ref{thm: exact berge}, the result holds by Theorem~\ref{thm: Kang Ni Shan},
and the case  $r=2$ in Theorem \ref{thm: exact berge} follows from the
result of Erd\H{o}s and Gallai \cite{gallai1959maximal}. %Tutte-Berge formula.
Thus, in the following, we may assume $s \le r \le 2s+1$ and $r \ge 3$.



\begin{proof}[\bf Proof of Theorem~\ref{thm: exact berge}] 
    Let us first consider the lower bound.
    %We first consider the lower bounds by several extremal hypergraphs satisfying conditions.
    For $r \le s+1$, we consider the following two $r$-graphs:
    one is an $r$-graph formed by the union of a complete $r$-graph of order \(2s+1\) and \(n-2s-1\) isolated vertices (i.e., vertices of degree 0); the other is an $n$-vertex $r$-graph where there exists a vertex subset $S$ of size $s$ such that every hyperedge intersects $S$ in at least $r-1$ vertices.
    For $s+2 \leq r \leq 2s$, the lower bound follows from the union of a complete $r$-graph of order \(2s+1\) and \(n-2s-1\) isolated vertices. For $r\geq2s+1$, the lower bound follows from any $r$-graph with $s$ hyperedges.
    It is easy to verify that these hypergraphs are $\mathcal{B}M_{s+1}$-free.
        %For $r \geq 2s+1$,  $n$ vertices and arbitrary $s$ hyperedges is the extremal graph.

    In the following, we prove the upper bound.
    Let $\mathcal{H}$ be a $\cB M_{s+1}$-free $r$-graph on $n$ vertices with maximum number of hyperedges.
    Note that adding a new hyperedge to $\cH$ increases the size of the core of a maximum Berge matching by at most $1$.
    Thus, by the maximality of $e(\cH)$, we may assume that $\mathcal{H}$ contains a $\cB M_s$, denoted by $\cM$. Let $\cE=\{e_1,\dots,e_s\}$ be the hyperedges of $\cM$, with the core $\cC = \{u_1,v_1,\dots,u_s,v_s\}$ such that $\{u_i,v_i\}\subseteq e_i$ for each $i\in [s]$.
    For a subset $S \subseteq \cC$, we write $\overline{S} = \{u_i \mid v_i \in S\} \cup \{v_i \mid u_i \in S\}$.
   % Note that for every $e \in E(\cH)\setminus \cE$, $|e\cap \cC|\geq s-1$, otherwise we can find a Berge copy of $M_{s+1}$ by adding the hyperedge $e$ to $\cE$.
    %For a vertex $v \notin \cC$, let $E_{\cC}(v) = \{e \in E(\cH)\setminus \cE \mid v \in e\}$, $d_{\cC}(v) = |E_{\cC}(v)|$.
    %Note that $d_{\cC}(v)$ is not exactly the degree of $v$ in $\cH$, since we do not count the hyperedges in $\cE$.
    %Let $N_\cC(v) = \{u \mid \exists e \in E_\cC(v), u \in e\}$ be the neighbors of $v$.
   % Note that $N_\cC(v) \subseteq \cC$, and $|N_\cC(v)| = r-1$ if and only if $d_\cC(v) = 1$.
   % Note that $|S| = |\overline{S}|$.
    % Set $d_{\cC}(v)=|N_{\cC}(v)|$. For every $v\in V(\cH)\setminus {\cC}$, let $C(v)=\{E\in N_{\cC}(v):~\{u_i,v_i\}\subseteq E \text{~for some~} i\in [k]\}$. And for every $v\in V(\cH)$, let $D(v)=\{u\in V(\cH):~\text{there exists $E\in E(\cH)$ such that $\{u,v\}\subseteq E$}\}$ be the neighbourhoods of $v$.

    Let $\cH'$ be the $r$-graph obtained from $\cH$ by deleting $s$ hyperedges of $\cE$. %Let $\mathcal{H}'=\mathcal{H}-\cE$.
    Let $E_{\cH'}(v)$ denote the set of hyperedges in $\cH'$ that contain $v$.
    Observe that each hyperedge of $\cH'$ intersects $\cC$ in at least $r-1$ vertices. Indeed, if there exists a hyperedge $e\in E(\cH')$ such that $e\setminus \cC=\{u,v\}$, then $\cE\cup \{e\}$ forms a $\cB M_{s+1}$ in $\cH$ with the core $\cC\cup\{u,v\}$, a contradiction.
    This implies that $N_{\cH'}(v) \subseteq \cC$ for any $v\in V(\cH)\setminus \cC$.
    %$u\in V(\cH)\setminus \cC$ such that $\{u,v\}\subseteq e\in E(\cH'then $\cE\cup \{e\}$ is a Berge copy of $M_{s+1}$  with the core $\cC\cup \{v,u\}$ of $\cH$, a contradiction.

    %The following claim holds for the case when $s\leq r\leq 2s$.
    \begin{cla}\label{claim: 2 isolated vtx}
        If $s\leq r\leq 2s$,
        then either $e(\cH)\leq \binom{2s+1}{r}$,
        or there are at least two vertices $u, v\in V(\cH)\setminus \cC$ such that $d_{\cH'}(u) \geq 1$ and $d_{\cH'}(v) \geq 1$.
    \end{cla}

    \begin{proof}[{\bf{Proof of Claim~\ref{claim: 2 isolated vtx}.}}]
    %\noindent\textbf{Proof of Claim~\ref{claim: 2 isolated vtx}:}
    It is sufficient to prove that if there is at most one vertex $u\in V(\cH)\setminus \cC$ with $d_{\cH'}(u) \geq 1$, then $e(\cH)\leq \binom{2s+1}{r}$.
    If $d_{\cH'}(v)=0$ for every $v\in V(\cH)\setminus \cC$, then $e(\cH) \leq s+\binom{2s}{r}\leq \binom{2s+1}{r}$ as $r\leq 2s$.

    Let $v\in V(\cH)\setminus \cC$ be the unique vertex with $d_{\cH'}(v)\geq 1$. If $e_i \subseteq \cC \cup \{v\}$ for every $i\in [s]$, then $e(\cH) \leq \binom{2s+1}{r}$ and we are done. Now suppose that there exists some $i\in [s]$ such that $e_i \nsubseteq \cC \cup \{v\}$.
    Without loss of generality, suppose $e_s$ contains a vertex $w \notin \cC \cup \{v\}$. Note that $|\cC|=2s$. If $r=2s$, then we have $u_s\in N_{\cH'}(v)$ or $v_s\in N_{\cH'}(v)$. Thus, there exists a $\cB M_{s+1}$ in $\cH$ by adding the hyperedge containing $v$ and $u_s$~(or $v_s$) to $\cE$, and extending the core to $\cC\cup \{v,w\}$, a contradiction. Now we assume that $s\leq r<2s$.
    %Note that $N_{\cH'}(v) \subseteq \cC$. Indeed, if there exists a vertex $u\in V(\cH)\setminus \cC$ such that $\{u,v\}\subseteq e\in E(\cH')$, then $\cE\cup \{e\}$ is a Berge copy of $M_{s+1}$  with the core $\cC\cup \{v,u\}$ of $\cH$, a contradiction.
    %Next we consider two cases: one is $e_i \subseteq \cC \cup \{v\}$ for every $i\in [s]$, and the other is that there exists $e_i \nsubseteq \cC \cup \{v\}$ for some $i\in [s]$.
    %For the first case, we have  $e(\cH) \leq \binom{2s+1}{r}$.
    We claim that $u_s\notin N_{\cH'}(v)$ and $v_s \notin N_{\cH'}(v)$.
    Indeed, if $u_s\in N_{\cH'}(v)$ (or $v_s\in N_{\cH'}(v)$), then as in the case $r=2s$ we may find a $\cB M_{s+1}$ in $\cH$, a contradiction.
    %If $r=2s$, then we have $u_s\in N_{\cH'}(v)$ or $v_s\in N_{\cH'}(v)$. So there exists a $\cB M_{s+1}$ in $\cH$, a contradiction.

    Then $|E_{\cH'}(v)|\leq \binom{|N_{\cH'}(v)|}{r-1}$.
    Note that $E(\cH'[\cC])=E(\cH')\setminus E_{\cH'}(v)$.
    %Note that every hyperedge of $E(\cH')\setminus E_{\cH'}(v)$ is contained in the vertex set $\cC$.
    We now partition $E(\cH'[\cC])$ into the following three classes:
    $\cE_1$ (consisting of hyperedges not containing $\overline{N_{\cH'}(v)}$),
    $\cE_2$ (consisting of hyperedges not containing $\{u_s,v_s\}$),
    and $\cE_3$ (consisting of hyperedges containing both a vertex of $\{u_s,v_s\}$ and a vertex of $\overline{N_{\cH'}(v)}$).
    %We now consider  the edges of $E(\cH')\setminus E_{\cH'}(v)$ formed by $\cE_1$ (hyperedges  not containing $\overline{N_{\cH'}(v)}$), $\cE_2$ (hyperedges not containing $\{u_s,v_s\}$),  and $\cE_3$ (hyperedges containing both  one vertex in $\{u_s,v_s\}$ and one  in $\overline{N_{\cH'}(v)}$).
    We assert that $\cE_3=\emptyset$. Otherwise, assume that $e\in \cE_3$ contains the vertex $u_s$ and $u_i\in \overline{N_{\cH'}(v)}$ (or $v_i\in \overline{N_{\cH'}(v)}$) for some $i\in [s-1]$.
    By the definition of $\overline{N_{\cH'}(v)}$, there is a hyperedge $e'\in E(\cH')$ containing $v$ and $v_i$ (or $u_i$). Then $(\cE\setminus \{e_i\})\cup \{e',e\}$ forms a $\cB M_{s+1}$ in $\cH$ with the core $\cC\cup \{v,w\}$, a contradiction.
    %i.e., for any hyperedge $e\in E(\cH')\setminus  E_{\cH'}(v)$, $e$ can not contain both one vertex in $\{u_s,v_s\}$ and one  in $\overline{N_{\cH'}(v)}$.
    %Suppose, without loss of generality, that $e\in \cE_3$ contains the vertex $u_s$ and $u_i\in \overline{N_{\cH'}(v)}$ for some $i\in [s-1]$.
    Note that $r-1\leq|N_{\cH'}(v)|\leq 2s-2$. Therefore, we have
     \begin{equation*}
    \begin{aligned}
        e(\cH)= s+|E_{\cH'}(v)|+|\cE_1|+|\cE_2| &\leq s + \binom{|N_{\cH'}(v)|}{r-1} + \binom{2s-|N_{{\cH'}}(v)|}{r} + \binom{2s-2}{r}
        %&\leq s + \binom{2s-2}{r-1} + \binom{2s-r+1}{r}+\binom{2s-2}{r} \\
        \le \binom{2s+1}{r}.
    \end{aligned}
     \end{equation*}
     \end{proof}
     %\hfill $\square$ \par

    %  Then $D(v)\subseteq U$. If for every $E_i$, $E_i\subseteq U$, then $|E(\cH)|\leq \binom{2k+1}{r}$.
    %  If there is some $E_i$, say $E_k$, with $E_k\setminus U\neq \emptyset$.
    %  Then for every $u_i\in D(v)$, there is no $E\in E(\cH)$ such that  $\{v_i,u_k\}\subseteq E$ or $\{v_i,v_k\}\subseteq E$. Otherwise, suppose $E'\in E(\cH)\setminus M$ with $\{v,u_i\}\subseteq E'$, then $(M\setminus E_i)\cup (E'\cup E_k)$ contains a copy of $\cB M_{k+1}$.
    %  Since $D(v)\geq k-1$, the number of hyperedges besides $M$ intersecting with $\{u_k,v_k\}$ is at most $\binom{|U\setminus D(v)|}{r}\leq k+1$.
    %  And the number of hyperedges besides $M$ do not intersect with $\{u_k,v_l\}$ is at most $\binom{|U\setminus\{u_k,v_k\}\cup \{v\}|}{k}=\binom{2k-1}{r}$. Thus, we have
    %  \begin{align*}
    %      |E(\cH)|\leq & |\{E\in E(\cH)\setminus M:~E\cap \{u_k,v_k\}\neq \emptyset\}|+|\{E\in E(\cH)\setminus M:~E\cap \{u_k,v_k\}= \emptyset\}|+|M|\\
    %      \leq & k+1+\binom{2k-1}{r}+k\leq \binom{2k+1}{r}
    %  \end{align*}
    %  When $k\leq r\leq 2k$.
    %  \hfill $\square$

    % \noindent\textbf{Situation (i): $r=s$.}
    \subsection{\normalsize Proof of Theorem~\ref{thm: exact berge} for $r=s$}

    \begin{cla}\label{claim: outdegree at least 2}
        If $d_{\cH'}(v)\leq 1$ for all $v\in V(\cH)\setminus \cC$,  then $e(\cH)\leq \max\left\{\binom{2s+1}{s},s(n-s)+1\right\}$.
    \end{cla}
    \begin{proof}[{\bf Proof of Claim \ref{claim: outdegree at least 2}.}]
    %\noindent\textbf{Proof of Claim~\ref{claim: outdegree at least 2}:}
    Note that $e(\cH'[\cC])\leq \binom{2s}{r}=\binom{2s}{s}$. Since $d_{\cH'}(v)\leq 1$ for all $v\in V(\cH)\setminus \cC$,
    \begin{align*}
        e(\cH) &= s+e(\cH'[\cC])+\sum_{v\in V(\cH)\setminus \cC}d_{\cH'}(v)\leq s+\binom{2s}{s}+n-2s=\binom{2s}{s}+n-s.
        %&\le \max \left\{\binom{2s+1}{s},1+s(n-s)\right\}~~~~(\text{since}~s\ge 3~\text{and}~n\ge 2s+2).
    \end{align*}
    If $e(\cH)>\max\left\{\binom{2s+1}{s},s(n-s)+1\right\}$, then we have $\binom{2s}{s}+n-s\geq \binom{2s+1}{s}+1$ and $\binom{2s}{s}+n-s\geq s(n-s)+2$. This implies that $\binom{2s}{s}\geq (s-1)(n-s)+2\geq (s-1)\big[\binom{2s}{s-1}+1\big]+2=(s-1)\binom{2s}{s-1}+s+1$, which is a contradiction when $s\geq2$.
    \end{proof}
    %\hfill $\square$

  % By Claim~\ref{claim: outdegree at least 2},  we may assume that  there exist some $v\in V(\cH) \setminus \cC$ with $d_{\cH'}(v) \ge 2$.

    \smallskip
    \noindent\textbf{Case 1.}  There exists $w \in V(\cH)\setminus \cC$ such that $\{u_i,v_i\} \subseteq N_{\cH'}(w)$ for some $i\in [s]$.
    \smallskip
    \smallskip

   %     In other words, one of the following holds:
   %     \begin{enumerate}
    %        \item  $\exists e_w \in E_{\cC}(w)$ such that $u_i,v_i \in e_w$.
    %        \item  $\exists e_{w1}, e_{w2} \in E_{\cC}(w)$ such that $u_i \in e_{w1}, v_i \in e_{w2}$.
    %    \end{enumerate}
    % We choose a such vertex $w$ such that $d_{\cH'}(w)$ achieves minimum.
    We choose such a vertex $w$ so that $d_{\cH'}(w)$ is minimal.
    Without loss of generality, assume that $\{u_1,v_1\} \subseteq N_{\cH'}(w)$. Then we have $\{u_1,v_1\} \subseteq e_w\in E_{\cH'}(w)$, or $u_1\in e_{w}^1$ and $v_1\in e_{w}^2$, where $e_{w}^1,e_{w}^2\in E_{\cH'}(w)$ and $e_{w}^1\neq e_{w}^2$.
        Let
        \[
        Y = N_{\cH'}\big(V(\cH)\setminus \left(\cC \cup\{w\}\right)\big).
        \]
    Clearly, $Y \subseteq \cC$. We claim that $u_1,v_1\notin Y$. Otherwise, there exists a $\cB M_{s+1}$ in $\cH$ with the core $\cC\cup\{w,z\}$, where $z\in V(\cH)\setminus \left(\cC \cup\{w\}\right)$ and $u_1\in N_{\cH'}(z)$ (or $v_1\in N_{\cH'}(z)$).

    By Claim~\ref{claim: 2 isolated vtx}, we may suppose that there exists at least one vertex in $V(\cH)\setminus \left(\cC \cup\{w\}\right)$ such that its degree in $\cH'$ is at least 1. Otherwise, $e(\cH)\leq \binom{2s+1}{r}$ and we are done.
    Thus, $|Y|\ge r-1=s-1$. %We next discuss two cases.

    \smallskip
    \noindent\textbf{Case 1.1.} $|Y| = s-1$.
    \smallskip
    \smallskip

    It is easy to see that $d_{\cH'}(v)\le 1$ for any $v\in V(\cH)\setminus \{\cC\cup w\}$. By Claim~\ref{claim: outdegree at least 2}, we may suppose $d_{\cH'}(w) \ge 2$. Otherwise, $e(\cH)\leq \max\left\{\binom{2s+1}{s},s(n-s)+1\right\}$ and we are done.
    Let $\cE_{1}$ and $\cE_{2}$ denote the set of hyperedges in $\cH'$ containing a vertex in $V(\cH)\setminus (\cC \cup \{w\})$, and the vertex $w$, respectively. Note that $\cE_{2}=E_{\cH'}(w)$. Let $\Omega = \{v_1,\ldots,v_s\}$ and $\cE_{3}=E(\cH'[\cC])\backslash E(\cH'[\Omega])$. Then
  \begin{align*}
  E(\cH)= \cE\cup E(\cH'[\Omega])\cup \bigcup_{i=1}^{3}\cE_{i}.
  \end{align*}
Clearly, $e(\cH'[\Omega])\le 1$. Note that $d_{\cH'}(v) \le 1$ for any $v \in V(\cH) \setminus (\cC \cup \{w\})$.
So $|\cE_{1}|\le n-2s-1$.
%It is sufficient to dic   the size of $\cE_2$ and $\cE_3$.



    Recall that $u_1,v_1\notin Y$. We claim that  $\{u_i,v_i\}\nsubseteq Y$ for all $2\leq i\leq s$. Otherwise, there exists some $j$ such that $\{u_j,v_j\}\subseteq Y$. Then $V(\cH) \setminus (\cC \cup \{w\})$  contains a vertex $v$ such that $\{u_j,v_j\} \subseteq N_{\cH'}(v)$ and $d_{\cH'}(v)\le 1$,  which contradicts the minimality of $d_{\cH'}(w)$.
    Without loss of generality, we assume that $Y = \{v_2,\ldots,v_s\}$.
    %Since $d_{\cH'}(w) \ge 2$, we have $|N_{\cH'}(w)| \ge r=s$.
    We claim that $N_{\cH'}(w) \subseteq \Omega \cup \{u_1\}$. Indeed, if $u_i\in N_{\cH'}(w)$ for some $2\le i \le s$, then there exists a $\cB M_{s+1}$, obtained by adding to $\cE\setminus \{e_i\}$ the hyperedge containing $wu_i$ and the hyperedge containing $v_iw'$ for some $w' \in V(\mathcal{H}) \setminus (\cC \cup \{w\})$ by the definition of $Y$.
    So $|\cE_2|\le \binom{s+1}{s-1}=\binom{s+1}{2}$.

    Since $d_{\cH'}(w) \ge 2$, we have $|N_{\cH'}(w)| \ge r=s$.
    Recall that $\{u_1,v_1\} \subseteq N_{\cH'}(w)$.
    Without loss of generality, we may assume that $\{u_1,v_1,\ldots,v_{s-1}\} \subseteq N_{\cH'}(w)$.
    For any hyperedge $e \in E(\cH'[\cC])$ containing $u_1$, $e$ does not contain any $u_j$ for $2\leq j\leq s$. Otherwise, by the definition of $Y$, we may find a $\cB M_{s+1}$ in $\cH$ whose core has edge set $\{wv_1,u_1u_j,v_jw',u_2v_2,\dots,u_sv_s\}\backslash\{u_jv_j\}$.
    For any hyperedge $e \in E(\cH'[\cC])$ containing $u_i$ with $2\leq i\leq s$, $e$ does not contain any other $u_j$ ($1\leq j\leq s$ and $j\neq i$) or $v_1$. Otherwise, by the definition of $Y$, we may find a $\cB M_{s+1}$ in $\cH$ (see Figure~\ref{fig: case 1.1}).
    It means that $|\cE_3|\le \binom{s}{r-1}+s-1=2s-1$.
    %For any hyperedge $e \in E(\cH'[\cC])$ containing $u_i$, $e$ does not contain any $u_j$ ($2\leq j\leq s$) or $v_1$. Otherwise, we may find a $\cB M_{s+1}$ in $\cH$~(see Figure~\ref{fig: case 1.1}). It means $|\cE_3|\le s$.

        \begin{figure}
            \centering
            \includegraphics[width=0.9\linewidth]{Berge-04.jpg}
            \caption{The illustration of $\cB M_{s+1}$ by the thick red lines. The left figure shows when a hyperedge $e \in \cE_3$ contains $u_i, u_j$~(the figure shows the case $2 \le i,j \le s-1$, the case when $i \in \{1,s\}$ holds similarly).
            The right figure shows when a hyperedge $e\in \cE_3$ contains $u_i, v_1$ ($2\leq i\leq s$).
            The hyperedge connecting $w'$ comes from the definition of $Y$.
            }
            \label{fig: case 1.1}
        \end{figure}

        Therefore, we have
        \begin{equation*}
        \begin{aligned}
            e(\cH)\leq   |\cE|+ e(\cH'[\Omega])+\sum_{i=1}^{3}|\cE_{i}| &\le s + 1+(n - 2s - 1) + \binom{s+1}{2} + 2s-1 \\
            & = n+s-1+\binom{s+1}{2} \\
            &\le \max\left\{ \binom{2s+1}{s}, s(n-s)+1 \right\},
        \end{aligned}
        \end{equation*}
        where the last inequality can be proved by a method similar to that of Claim \ref{claim: outdegree at least 2} when $s\geq3$.

    \smallskip
    \noindent\textbf{Case 1.2.} $|Y| \ge s$.
    \smallskip
    \smallskip
        %Now we consider $e_1 \in \cE$ containing $u_1v_1$.

    Let $\cE_{1}$ and $\cE_{2}$ denote the set of  hyperedges in $\cH'$ containing the vertex $w$, and a vertex in $V(\cH)\setminus (\cC \cup \{w\})$, respectively. Let $\cE_{3}=E(\cH'[\cC])$.
    %and the some vertices in $\cC$, and
    Then
    \begin{align*}
    E(\cH)= \cE\cup \bigcup_{i=1}^{3}\cE_{i}.
    \end{align*}

    Recall that $\{u_1,v_1\}\subseteq N_{\cH'}(w)$. %and
    We claim that $e_1 \subseteq \cC\cup \{w\}$. Otherwise, we may find a $\cB M_{s+1}$ in $\cH$ with the core $\cC\cup\{w,z\}$, where $z\in e_1\setminus \left(\cC\cup \{w\}\right)$.
    So we can choose a set $X$ of $s$ vertices from $(e_1 \cup e_w) \cap \cC$ or $(e_1 \cup e_{w}^1 \cup e_{w}^2) \cap \cC$ such that  $\{u_1,v_1\}\subseteq X$.

    Recall that $u_1,v_1\notin Y$. It follows that $u_1,v_1\notin \overline{Y}$ and $|\overline{X} \cap \overline{Y}| \le s-2$.
    We claim that $\overline{X} \cap \overline{Y} = \emptyset$.
    Otherwise, assume that $u_j\in \overline{X} \cap \overline{Y}$ (or $v_j\in \overline{X} \cap \overline{Y}$)  for some $2 \le j \le s$.
    Then for the hyperedge $e_j \in \cE$, apart from $u_j$~(or $v_j$), it does not contain vertices of $V(\cH)\setminus \cC$, vertices in $\overline{X}$, and vertices in $\overline{Y}$. Otherwise, we may obtain a $\cB M_{s+1}$ in $\cH$~(see Figure~\ref{fig: Berge exact} for illustration), a contradiction.
    Thus, we have
        \begin{equation*}
        \begin{aligned}
            s- 1=|e_j| -1 \le |\cC \setminus (\overline{X} \cup \overline{Y})| = 2s-|X|-|Y| + |\overline{X} \cap \overline{Y}| \le s - |Y| + s - 2.
        \end{aligned}
        \end{equation*}
        It yields $|Y| \le s-1$, which contradicts $|Y|\ge s$. Thus, we have $\overline{X} \cap \overline{Y} = \emptyset$.

        \begin{figure}[t]
            \centering
            \includegraphics[width=0.9\textwidth]{Berge-03.jpg}
            \caption{The four cases: $e_j$ contains $w$, $w'' \neq w$, a vertex in $\overline{X}$, or a vertex in $\overline{Y}$. In each case, we can find a $\cB M_{s+1}$ using either red lines or dotted lines.
            In the picture, we only list some typical cases. For example, when $e_j$ contains $u_2 \in \overline{X}$, we only consider $u_2 \in e_j\cap \overline{X}$.
            However, other cases hold similarly.}\label{fig: Berge exact}
        \end{figure}

        Combining $|X|=|\overline{X}| = s$, $|Y|=|\overline{Y}| \ge s$, $|\overline{X} \cup \overline{Y}| \le |\cC|= 2s$, and  $|\overline{X} \cap \overline{Y}| = 0$, we have $|Y| = |\overline{Y}| = s$ and $|X \cap Y |= 0$.
        We claim that $X \cap \overline{Y} = \emptyset$. Otherwise, there exists a $\cB M_{s+1}$ in $\cH$. So we have $X = \overline{X}$ and $Y = \overline{Y}$.
        Thus, all hyperedges of $\cE_1$ are contained in $\{w\} \cup X$ (otherwise, there exists a $\cB M_{s+1}$ in $\cH$ since $Y = \overline{Y}$).
        Then $|\cE_1|\leq \binom{|X|}{s-1}=\binom{s}{s-1}$.
        Furthermore, if $e\in\cE_3$ contains a vertex from $X$, then $e=X$, and if $e\in\cE_3$ contains a vertex from $Y$, then $e=Y$ (otherwise, there exists a $\cB M_{s+1}$ in $\cH$).
        Then $|\cE_3|\leq 2$.
        By the definition of $Y$, we have $|\cE_2|\leq \binom{|Y|}{s-1}(n-2s-1)=\binom{s}{s-1}(n-2s-1)$.
        %all hyperedges of $\cE_2$ are contained in $\{w'\}\cup Y$, and the hyperedges of $\cE_3$ either contains $u_i \in X$ or $v_i \in X$, or contains $u_i \in Y$ or $v_i \in Y$.
        Therefore,
        \begin{equation*}
        \begin{aligned}
            e(\cH) &= |\cE|+\sum_{i=1}^{3}|\cE_{i}|\\
            & \leq s + \binom{s}{s-1}  +\binom{s}{s-1}(n - 2s - 1) + 2 \\
            & = s(n - s) + 1 - (s^2 -s -1) \le s(n-s)+1.
        \end{aligned}
        \end{equation*}


    \smallskip
    \noindent\textbf{Case 2.} For any $v \in V(\cH)\setminus \cC$, $\{u_i,v_i\} \nsubseteq N_{\cH'}(v)$ for each $i\in [s]$.
    \smallskip
    \smallskip

    By Claim~\ref{claim: outdegree at least 2}, we may suppose $d_{\cH'}(w)\ge 2$ for some $w\in V(\cH)\setminus \cC$. Then we have that $N_{\cH'}(w) \subseteq \cC$ and $|N_{\cH'}(w)| \ge r=s$.
    Note that $\{u_i,v_i\}\nsubseteq N_{\cH'}(w)$ for every $i\in[s]$.
    Hence, without loss of generality, we may assume that $N_{\cH'}(w) = \Omega$.

    By Claim~\ref{claim: 2 isolated vtx}, we may suppose that there exists at least one vertex in $V(\cH)\setminus \left(\cC \cup\{w\}\right)$ such that its degree in $\cH'$ is at least 1. Otherwise, $e(\cH)\leq \binom{2s+1}{r}$ and we are done.
    It follows that $N_{\cH'}(w') \subseteq \Omega$ for any $w' \in V(\cH)\setminus \left(\cC \cup\{w\}\right)$. Otherwise, there exists a $\cB M_{s+1}$ in $\cH$ with the core $\cC\cup\{w,w'\}$.

    % No edge in $E(\cH)\setminus \cE$ contains $w \in V(\cH)\setminus \cC$ and $u_i,v_i$ for some $i$.

    % By the special properties of $v_i$, $N_{\cC}(w) \subseteq \Omega = \{v_1,\ldots,v_s\}$ for all $w\in V(\cH)\setminus\cC$.

    % \begin{cla}\label{claim: cover once}
    %     If there exists a set $W \subseteq V(\cH)\setminus \cC$ such that each $v_i$, $i=1,2,\ldots,s$ is in at least one set in ${\{N_{\cC}(w)\}}_{w\in W}$.
    % Then every $e_i$ can not contain any vertex outside $\cC \cup W$.
    % Also, if $e$ is an edge from $E(\cH) \setminus \cE$, if $e$ contains some $u_i$, then $e \subseteq \cC \cup W$.
    % \end{cla}


    \smallskip
    \noindent\textbf{Case 2.1.}
    There exists a set $W \subseteq V(\cH)\setminus \cC$ such that each $v_i$ $(i\in [s])$ is contained in at least two members of the family $\{N_{\cH'}(v)\mid	v\in W\}$.
    \smallskip
    \smallskip

    It follows that for each $e \in E(\mathcal{H})$, the hyperedge $e$ can not contain both $u_i$ and $u_j$ for any $i \neq j$.
    Otherwise, we would get a $\cB M_{s+1}$ in $\cH$ by connecting $u_iu_j$ with the edge, and connecting $u_i,u_j$ to two different elements in $W$ by definition.
    Then every hyperedge $e \in E(\mathcal{H})$ intersects $\Omega$ in at least $s-1$ vertices.
    % Similarly, for any hyperedge $e\in E(\cH')$, $e$ intersects $\Omega$ in at least $s-1$ vertices.
    Therefore, we have
    \begin{equation*}
    \begin{aligned}
        e(\cH) & \leq 1 + \binom{s}{s-1}(n-s) = 1 +s(n-s).
    \end{aligned}
    \end{equation*}


    \smallskip
    \noindent\textbf{Case 2.2.} There does not exist $W \subseteq V(\cH)\setminus \cC$ such that each $v_i$ $(i\in [s])$ is contained in at least two members of the family $\{N_{\cH'}(v)\mid v\in W\}$. %For any $W \subseteq V(\cH)\setminus \cC$, there exist some $v_i$ are  in at most  one set of the family $\{N_{\cH'}(v)\mid	v\in W\}$, where  $i\in [s]$.
    \smallskip
    \smallskip

     % Let $\cE_{1}$,  $\cE_{2}$,  $ \cE_{3}$, and $ \cE_{4}$ denote the set of  hyperedges   containing  $w$, $w' \in V(\cH)\setminus (\cC \cup \{w\})$, $u_1$, and  some vertices in $\cC\setminus \{u_1\}$, respectively.
%  Then
 % \begin{align*}
 % E(\cH)= \cE\cup \bigcup_{i=1}^{4}\cE_{i}.
 % \end{align*}

    %Recall that there exists at least one vertex in $V(\cH)\setminus \left(\cC \cup\{w\}\right)$ such that its degree in $\cH'$ is at least 1.
    Note that $N_{\cH'}(w) = \Omega$ and $N_{\cH'}(w') \subseteq \Omega$ for any $w' \in V(\cH)\setminus (\cC \cup \{w\})$. It follows that for every $w' \in V(\cH)\setminus (\cC \cup \{w\})$, we have $d_{\cH'}(w') \le 1$.
    Otherwise, there exists a set $W=\{w,w'\}$ that satisfies Case~2.1, a contradiction.
    So, by Claim \ref{claim: 2 isolated vtx}, there exists one vertex $w'\in V(\cH)\setminus \left(\cC \cup\{w\}\right)$ such that $d_{\cH'}(w')=1$.
    Similarly, if $w_1',w_2' \in V(\cH)\setminus (\cC \cup \{w\})$ satisfies $d_{\cH'}(w_1') = d_{\cH'}(w_2') = 1$,
    then $N_{\cH'}(w_1') = N_{\cH'}(w_2')$.
    %Otherwise, we reduce to Case 2.1.
    Without loss of generality, assume that $N_{\cH'}(w') = \Omega\setminus \{v_1\}$ for all $w' \in V(\cH)\setminus (\cC \cup \{w\})$ with $d_{\cH'}(w') = 1$.
        Let $\cE_{1}$ and $\cE_{2}$ denote the set of  hyperedges in $\cH'$ containing the vertex $w$, and a vertex in $V(\cH)\setminus (\cC \cup \{w\})$, respectively. Let $\cE_{3}=E(\cH'[\cC])$.
Clearly, we have $|\cE_1| \le \binom{s}{s-1} = s $, $|\cE_2| \le n-2s-1$.

Now we consider the hyperedges in $\cE_{3}$.
If $e \in \cE_{3}$ contains $u_i$~($1 \le i \le s$), then it must lie in $\{u_i\}\cup \Omega$, at most $\binom{s}{s-1} = s$ such hyperedges. Otherwise, we would get a $\cB M_{s+1}$ in $\cH$ by connecting $u_iu_j$ with the edge, and connecting $u_i,u_j$ to $w,w'$.
%If $e \in \cE_{3}$ contains $u_i$~($2 \le i \le s$), then it must be the hyperedge $\{u_i\} \cup \Omega \setminus \{v_1\}$, at most $s-1$ such hyperedges.
With only one hyperedge on $\Omega$, we have $|\cE_3| \le s^2+1$.

    % Therefore, the number of hyperedges in $\cH'$ containing a vertex in $V(\cH)\setminus (\cC \cup \{w\})$ is at most $n-2s-1$.

    % the number of hyperedges in $E(\cH'[\cC])$ containing $u_1$...

    % the number of hyperedges in $E(\cH'[\cC])$ containing $u_i$...


    % Note that the number of hyperedges in $\cH'$ containing $w$ is at most $\binom{s}{s-1}$. Thus,


    %some $w' \in V(\cH)\setminus (\cC \cup \{w\})$, $u_1$, and some $u_i$ ($2 \le i \le s$)  are at most $\binom{s}{s-1}$,  $n - 2s - 1$, $\binom{s+1}{s-1}$, and $s-1$, respectively.
    This implies that
    \begin{equation*}
    \begin{aligned}
        e(\cH) & \le |\cE|+|\cE_1| + |\cE_2| + |\cE_3|\\
        % e(\cH) & \le |\cE|+e(\cH[\Omega])+\binom{s}{s-1}+(n - 2s - 1)+\binom{s+1}{s-1}+(s-1)\\
        &\leq s + s + n-2s-1 +s^2+1 \\
               & = n + s^2  \le \max\left\{ \binom{2s+1}{s}, 1+s(n-s) \right\},
    \end{aligned}
    \end{equation*}
    where the last inequality can be proved by a method similar to that of Claim \ref{claim: outdegree at least 2} when $s\geq2$.

    % then by Claim~\ref{claim: outdegree at least 2}, and there are two vertices $w_1,w_2 \in V(\cH)\setminus \cC$ with $d_{\cC}(w_i) \geq 1$ for $i=1,2$ and $N_{\cC}(w_1) \neq N_{\cC}(w_2)$.
    % Without loss of generality, we may assume $e_1' = \{v_1,v_2,\ldots,v_{s-1}, w_1\}$ is an edge and $e_2' = \{v_2,\ldots,v_s, w_2\}$ is an edge.

    % Clearly, $e_1\setminus\{u_1\} \subseteq \{v_1,\ldots,v_s, w_1\}$ and $e_s\setminus\{u_s\} \subseteq \{v_1,\ldots,v_s, w_2\}$.
    % For $2 \le i \le s-1$, $e_i\setminus \{u_i\} \subseteq \{v_1,\ldots,v_s\}$.
    % Then $E(\cH)$ consists of
    % \begin{enumerate}
    %     \item $\cE$, $s$ hyperedges;
    %     \item $e$ containing some $u_i$, then $e\setminus \{u_i\} \subseteq \{v_1,\ldots,v_s\}$, at most $s \binom{s}{s-1}$ hyperedges;
    %     \item $e$ containing some $w \in V(\cH)\setminus \cC $, at most $n-2s$ hyperedges.
    % \end{enumerate}
    % In total,
    % \begin{equation*}
    % \begin{aligned}
    %     e(\cH) & \leq s + s\binom{s}{s-1} + (n-2s) \\
    %            & = 1 + s(n-s).
    % \end{aligned}
    % \end{equation*}

    % \textbf{Case 2.4:}
    % The only last case is that for every $v\in V(\cH)\setminus \cC$, $d_{\cC}(v)\leq 1$.
    % And we have
    % \begin{equation*}
    % \begin{aligned}
    %     e(\cH) & \leq s + \binom{2s}{s} + (n - 2s) \leq \max\{\binom{2s+1}{s}, 1 + s(n-s)\}.
    % \end{aligned}
    % \end{equation*}


    \subsection{\normalsize Proof of Theorem~\ref{thm: exact berge} for $s+1 \le r \le 2s+1$}
    % \vspace*{.4cm}

    \smallskip
    \noindent\textbf{Case 1.} $r=s+1$.
    \smallskip
    \smallskip

    By Claim~\ref{claim: 2 isolated vtx}, if $e(\cH)\leq \binom{2s+1}{r}$, then the result holds.
    So we assume that $w_1,w_2$ are two vertices  such that $d_{\cH'}(w_i)\geq 1$ for $i=1,2$.
    Then $|N_{\cH'}(w_i)|\geq r-1=s$.
    And $\overline{N_{\cH'}(w_1)}\cap N_{\cH'}(w_2)=\emptyset$, otherwise we can find a Berge copy of $M_{s+1}$. Thus, $|N_{\cH'}(w_1)| = |N_{\cH'}(w_2)| = s$.

    For the case $N_{\cH'}(w_1)\cap N_{\cH'}(w_2)=\emptyset$, we write $N_{\cH'}(w_1)=V_1$ and $N_{\cH'}(w_2)=V_2$. Then for every $x\in V_1$ and $y\in V_2$, there is no hyperedge $e \in E(\cH)\setminus \cE$ with $\{x,y\}\subseteq e$, otherwise we can find a $\cB M_{s+1}$. Therefore, $E(\cH'[\cC])=\empty$ and for every vertex $v\in V(\cH)\setminus \cC$, either $N_{\cH'}(v) \subseteq V_1$ or $N_{\cH'}(v) \subseteq V_2 $. Thus, $e(\cH)\leq s+n-2s = n-s$.


    For the case $N_{\cH'}(w_1)\cap N_{\cH'}(w_2) \neq \emptyset$,  there exists $u_j$~(or $v_j$) in $\overline{N_{\cH'}(w_1)}\cap \overline{N_{\cH'}(w_2)}$. Now we consider the hyperedge $e_j$ in $\cE$.
    Apart from $u_j$~(or $v_j$), $e_j$ does not contain vertices in $V(\cH)\setminus \cC$,  $\overline{N_{\cH'}(w_1)}$, and  $\overline{N_{\cH'}(w_2)}$ by a similar argument as Figure~\ref{fig: Berge exact}.
    Then we have
    \begin{align*}
    s = |e_j| -1 \le  |\cC \setminus (\overline{N_{\cH'}(w_1)} \cup \overline{N_{\cH'}(w_2)})| \le |\overline{N_{\cH'}(w_1)} \cap \overline{N_{\cH'}(w_2)}|.
    \end{align*}
     It follows that $\overline{N_{\cH'}(w_1)} = \overline{N_{\cH'}(w_2)}$ and thus $N_{\cH'}(w_1) = N_{\cH'}(w_2)$.
    This only happens when $N_{\cH'}(w_1)$ contains exactly one vertex from $\{u_i,v_i\}$.
    Without loss of generality, we may assume $N_{\cH'}(w_1) = N_{\cH'}(w_2) = \Omega=\{v_1,\dots,v_s\}$.
    In this case, for every $w \in V(\cH) \setminus \cC$, either $N_{\cH'}(w) = \Omega$ or $d_{\cH'}(w) = 0$.
    In addition, every hyperedge containing $u_i$~(including $e_i$) contains $\Omega$.
    Thus, we have
    \[
        e(\cH)\leq e(\cH[\Omega])+\sum_{u\in V(\cH)\setminus\Omega }d_{\cH'}(u)\leq \binom{s}{s+1}+(n-s)=n-s.
    \]


    \smallskip
    %\vspace*{.4cm}
    \noindent\textbf{Case 2.} $s+2\leq r\leq 2s$.
    \smallskip
    \smallskip

    Suppose that the result does not hold. By Claim~\ref{claim: 2 isolated vtx}, we may assume $w_1,w_2\in V(\cH)\setminus \cC$ with $d_{\cH'}(w_i)\geq 1$ for $i=1,2$. Then $N_{\cH'}(w_i)\geq r-1\geq s+1$. It implies $\overline{N_{\cH'}(w_1)}\cap N_{\cH'}(w_2)\neq \emptyset$. By a similar argument, we can find a $\cB M_{s+1}$, a contradiction.


    \smallskip
    %\vspace*{.4cm}
    \noindent\textbf{Case 3.} $r = 2s+1$.
    \smallskip
    \smallskip

    Suppose that the result does not hold, i.e., $e(\cH)\geq s+1$. Then there is at least one hyperedge $e \in E(\cH)\setminus \cE$. Since $r = 2s+1$, $e$ must be of the form $\{w\} \cup \cC$ for a $w \in V(\mathcal{H}) \setminus \cC$.
    Then $e_1$ contains a vertex outside $\{w\} \cup \cC$, and thus we may find a $\cB M_{s+1}$, a contradiction.

   % Suppose otherwise $e(\cH)\geq s+1$, then there are at least two hyperedges $e_1',e_2'\in E(\cH)\setminus \cE$. If $|e_i' \setminus \cC|\geq 2$, then we can find a Berge copy of $M_{s+1}$, for $i=1,2$ Thus, $\cC \subseteq e_i'$ for $i=1,2$ and let $x_i=e_i'\setminus \cC$ for $i=1,2$.
   % As a result, $(\cM\setminus e_1)\cup (e'_1\cup e_2')$ contains a Berge copy of $M_{s+1}$ with core $\cC \cup\{x_1,x_2\}$. A contradiction.
\end{proof}






% % -------------------------------------------
% \section{\normalsize Proof of Theorems~\ref{thm: berge matching and a hypergraph}}\label{sec: berge matching and a hypergraph}
% %Let $\cH$ be an $r$-graph and $U\subseteq V(\cH)$ be a nonempty subset. Let $\cH[U]$ denote the subhypergraph of $\cH$ induced by $U$.
% Given a hypergraph $\cH$ and a vertex $v\in V(\cH)$, the link hypergraph $\cL_v$ is defined as
% \begin{eqnarray*}
% \cL_v=\big\{e\backslash\{v\}\,\big|\,e\in E(\cH), v\in e\big\}.
% \end{eqnarray*}

% Let us start with the proof of Theorem~\ref{thm: berge matching and a hypergraph} that we restate here for convenience.

% \begin{thm*}
% Let integers $3\leq r\leq s+1$ and $n$ be sufficiently large. Suppose $\cF$ is an $r$-graph with $\chi(\cF)>2$, or $\chi(\cF)=2$ and $q(\cF)=\infty$. %$F$ is a nonempty $r$-graph such that deleting any strongly independent set leaves it nonempty.
% Then
% \[
%     \ex_r(n,\mathcal{B}M_{s+1}\cup\{\cF\})=\ex_r(s,\mathcal{D}(\cF))+ \binom{s}{r-1}(n-s),
% \]
% where $\mathcal{D}(\cF)$ is the family of $r$-graphs obtained by deleting a strongly independent set from $\cF$.
% \end{thm*}

% \begin{proof}[\bf Proof] %of Theorem~\ref{thm: berge matching and a hypergraph}.]
%     %For a $r$-graph $\cF$ with $\chi(\cF) > 2$, or $\chi(\cF) = 2$ and $q(\cF) = \infty$.
%     Observe that for any strongly independent set $S$ of $\cF$, the $r$-graph $\cF - S$ is not empty. Otherwise, $\chi(\cF) = 2$ and $q(\cF)< \infty$, a contradiction.
%     This implies that $\mathcal{D}(\cF)$ does not contain the empty graph.
%     Thus there exists an $s$-vertex $\mathcal{D}(\cF)$-free $r$-graph $\cH_0$ with $\ex_r(s,\mathcal{D}(\cF))$ hyperedges.
%     % Observe that a nonempty $r$-graph $\cH$ contains a strongly independent set such that deleting it results in the empty graph if and only if $\chi(\cH)=2$ and $q(\cH)<\infty$. Therefore, a nonempty $r$-graph $\cH$ has the property that deleting any strongly independent set leaves it nonempty if and only if $\chi(\cH)>2$, or $\chi(\cH)=2$ and $q(\cH)=\infty$. This implies that there are no empty graphs in $\mathcal{F}$.
%     %Let $\cH$ be an $s$-vertex $\mathcal{D}(\cF)$-free $r$-graph with $\ex_r(s,\mathcal{D}(\cF))$ edges.
%     Now let us add $n-s$ new vertices and add all hyperedges containing one new vertex and $r-1$ vertices in $\cH_0$.
%     The resulting $r$-graph is $\mathcal{B}M_{s+1}$-free since $|V(\cH_0)|<s+1$, and $\cF$-free since $\cH_0$ is $\mathcal{D}(\cF)$-free and the $n-s$ vertices outside $\cH_0$ form a strongly independent set.
%     The construction gives the lower bound.

%     Now we consider the upper bound. Let $\cH$ be an $n$-vertex $\mathcal{B}M_{s+1}\cup\{\cF\}$-free $r$-graph with $e(\cH) > \ex_r(s, \mathcal{D}(\cF)) + \binom{s}{r-1}(n-s)$.
%     Let $V(\cH)=\{v_1,v_2,\dots,v_n\}$ and $d_\cH(v_1)\geq\cdots\geq d_\cH(v_n)$.

%     \begin{cla}\label{clam1}
%     $d_\cH(v_{s+1})\leq s+\binom{2s}{r-1}$.
%     \end{cla}
%     \begin{proof}[{\bf{Proof of Claim~\ref{clam1}.}}]
%     Suppose otherwise that $d_\cH(v_{s+1})> s+\binom{2s}{r-1}$. We may greedily find a $\mathcal{B}M_{s+1}$ in the following way. First, we pick a hyperedge $e_1$ containing $v_1$ and a vertex $v'_1$ distinct from $v_1,v_2,\dots,v_{s+1}$, and take $v_1$ and $v'_1$ as the core of $e_1$. In step $i$ ($2\leq i\leq s+1$), we pick a hyperedge $e_i$ containing $v_i$ and a vertex $v'_i$ distinct from $v_1,v_2,\dots,v_{s+1},v'_1,\dots,v'_{i-1}$ which is different from all previously selected hyperedges, and take $v_i$ and $v'_i$ as the core of $e_i$.
%     This is possible since the number of the previously selected hyperedges is at most $s$, and there are at most $\binom{s+i-1}{r-1}\leq \binom{2s}{r-1}$ hyperedges containing $v_i$ and $r-1$ vertices from $\{v_1,\dots,v_{s+1},v'_1,\dots,v'_{i-1}\}\backslash v_i$.
%     Then $e_1,\dots,e_{s+1}$ forms a $\mathcal{B}M_{s+1}$ whose core is $\{v_1v'_1,\dots,v_{s+1}v'_{s+1}\}$, a contradiction.
%     \end{proof}

%     Let $\cM_t$ be a maximum Berge matching in $\cH$, and let $A$ be the core of $\cM_t$. Then $|A| = 2t \le 2s$. Let ${\{e_i\}}_{1 \le i \le t}$ be the hyperedges of $\cM_t$.
%     Every hyperedge in $E(\cH)$ other than ${\{e_i\}}_{1 \le i \le t}$ intersects $A$ in at least $r-1$ vertices, otherwise a larger Berge matching would appear.
%     Assume $A=\{a_1,a_2,\ldots,a_{2t}\}$.
%     This implies that the number of hyperedges in $\cH$ is at most
%     \begin{eqnarray}
%     t+\frac{1}{r-1}\sum_{i=1}^{2t}d_\cH(a_i)\le t+\frac{1}{r-1}\sum_{i=1}^{2s}d_\cH(v_i)\le s+\frac{1}{r-1}\sum_{i=1}^{s}d_\cH(v_i)+\frac{s}{r-1}\bigg[s+\binom{2s}{r-1}\bigg].
%     \end{eqnarray}
%     The upper bound in this inequality follows from Claim~\ref{clam1}.

%     \begin{cla}\label{clam2}
%     $d_\cH(v_{i}) \le \binom{s-1}{r-2}(n-1) + \binom{s-1}{r-1}$.
%     \end{cla}

%     \begin{proof}[{\bf{Proof of Claim~\ref{clam2}.}}]
%      Let $d:=d(s,r)$ be a large enough fixed constant. If $d_\cH(v_i)\leq d$, then $d_\cH(v_i)\leq \binom{s-1}{r-2}(n-1)+\binom{s-1}{r-1}$ as $\binom{s-1}{r-2}\geq1$ and $n$ is sufficiently large. If $d_\cH(v_i)>d$, then we claim that the $(r-1)$-graph $\cL_{v_i}$ is $(r-1)$-uniform $\cB M_s$-free.
%         Suppose otherwise that $\cL_{v_i}$ contains a $\cB M_s$, denoted by $\cM_s$. Let $E(\cM_s)=\{e'_1,\dots,e'_s\}$.
%         Since $n$ is sufficiently large, we have $e(\cL_{v_i})=d_\cH(v_i)>d>s+\binom{2s}{r-1}$.
%         Then there exists an $(r-1)$-edge $e'\in E(\cL_{v_i})$ with $e'\neq e'_1,\dots,e'_s$ that contains a vertex $v'$ outside the core of $\cM_s$. Then $e'_1\cup\{v_i\},\dots,e'_s\cup\{v_i\},e'\cup\{v_i\}$ forms a $r$-uniform $\mathcal{B}M_{s+1}$ whose core is the core of $\cM_s$ together with the vertices $v_i$ and $v'$, a contradiction. By Theorem~\ref{thm: Khormali Palmer}, we have $e(\cL_{v_i})=d_\cH(v_i)\leq \binom{s-1}{r-2}(n-1)+\binom{s-1}{r-1}$.
%     \end{proof}

%     \begin{cla}\label{clam3}
%     $d_\cH(v_{s})> \binom{s-1}{r-2}n-c$ for some $c=c(s,r)>0$.
%     \end{cla}

%     \begin{proof}[{\bf{Proof of Claim~\ref{clam3}.}}]
%     Suppose to the contrary that $d_\cH(v_{s})\leq\binom{s-1}{r-2}n-c$ for any $c$. By Claim~\ref{clam2}, we have
%     % Since $n$ is sufficiently large, there exists some $c_0$ such that $$\frac{s-1}{r-1}\bigg[\binom{s-1}{r-2}(n-1)+\binom{s-1}{r-1}\bigg]+\binom{s-1}{r-2}\frac{n}{r-1}-\frac{c_0}{r-1}+\frac{s}{r-1}\bigg[s+\binom{2s}{r-1}\bigg]+s\le \binom{s}{r-1}(n-s).$$

%     % Let $d$ be a large enough fixed constant. For $1\leq i\leq s-1$, if $d_\cH(v_i)\leq d$, then $d_\cH(v_i)\leq \binom{s-1}{r-2}(n-1)+\binom{s-1}{r-1}$ as $\binom{s-1}{r-2}\geq1$ and $n$ is sufficiently large. If $d_\cH(v_i)>d$, then we claim that $\cL_{v_i}$ is $(r-1)$-uniform $\mathcal{B}M_s$-free. Therefore, $d_\cH(v_i)\leq \binom{s-1}{r-2}(n-1)+\binom{s-1}{r-1}$ for every $1\leq i\leq s-1$. By (3.1),
%     \begin{eqnarray*}
%     e(\cH)&\leq& \frac{1}{r-1}\sum_{i=1}^{s}d_\cH(v_i)+\frac{s}{r-1}\bigg[s+\binom{2s}{r-1}\bigg]+s \\
%     &\leq& \frac{s-1}{r-1}\bigg[\binom{s-1}{r-2}(n-1)+\binom{s-1}{r-1}\bigg]+\binom{s-1}{r-2}\frac{n}{r-1}-\frac{c}{r-1}+\frac{s}{r-1}\bigg[s+\binom{2s}{r-1}\bigg]+s \\
%     &=& \binom{s}{r-1}(n-s) + C - \frac{c}{r-1},
%     \end{eqnarray*}
%     where $C$ is a constant depending only on $r$ and $s$. When $c$ is large, we have $e(\cH)\leq \binom{s}{r-1}(n-s)$, a contradiction.
%     \end{proof}

%     Let $U=\{v_1,\dots,v_s\}$. If there exists a hyperedge $e$ intersecting $V(\cH)\backslash U$ in at least two vertices, then we take $u,u'\in e\backslash U$ as the core of $e$. Since $d_\cH(v_1)\geq\cdots\geq d_\cH(v_{s})> \binom{s-1}{r-2}n-c$ by Claim~\ref{clam3}, we may greedily find a $\mathcal{B}M_{s+1}$ in a way analogous to the proof of Claim~\ref{clam1}. Thus, $\cH$ contains only two types of hyperedges: those containing exactly $r-1$ vertices from $U$, and those containing exactly $r$ vertices from $U$. The number of hyperedges of the first type is at most $\binom{s}{r-1}(n-s)$, and the number of hyperedges of the second type is equal to $e(\cH[U])\leq \binom{s}{r}$.

%     If $\cH[U]$ is $\cD(\cF)$-free, then $e(\cH)\leq \ex_r(s,\cD(\cF))+\binom{s}{r-1}(n-s)$ and we are done.
%     Assume now that $\cH[U]$ contains an $r$-graph $\cF_0\in\cD(\cF)$.
%     Then for each $(|V(\cF)|-|V(\cF_0)|)$-set $X$ outside $U$, this copy of $\cF_0$ could be extended to a copy of $\cF$ in $V(\cF_0)\cup X$ using only hyperedges that contains exactly $r-1$ vertices from $V(\cF_0)$ and one vertex from $X$, by the definition of $\cD(\cF)$.
%     This means that at least one such hyperedge is missing from $\cH$. Let us partition the $n-s$ vertices outside $U$ to $\lfloor (n-s)/(|V(\cF)|-|V(\cF_0)|)\rfloor$ sets of size $|V(\cF)|-|V(\cF_0)|$ and a smaller set, arbitrarily.
%     Then we have a distinct missing hyperedge for each of the sets of size $|V(\cF)|-|V(\cF_0)|$, thus there are at most $\binom{s}{r-1}(n-s)-\lfloor (n-s)/(|V(\cF)|-|V(\cF_0)|)\rfloor$ hyperedges containing exactly $r-1$ vertices from $U$.
%     Therefore, the number of hyperedges of $H$ is at most $\binom{s}{r}+\binom{s}{r-1}(n-s)-\lfloor (n-s)/(|V(\cF)|-|V(\cF_0)|)\rfloor$.
%     Since $n$ is sufficiently large, we have $\binom{s}{r}-\lfloor (n-s)/(|V(F)|-|V(\cF_0)|)\rfloor\leq \ex_r(s,\cD(\cF)) + \binom{s}{r-1}(n-s)$. This completes the proof.
% \end{proof}

% \section{\normalsize Proofs of Theorems~\ref{thm: berge matching and berge bigraph}, \ref{thm: berge matching and berge Kpq} and \ref{prop2.8}}\label{sec: berge matching and berge bigraph}

% Let $S_\ell$ be a star with $\ell$ edges. The non-leaf vertex is called the \textit{center} of $S_\ell$.
% Khormali and Palmer~\cite{KHORMALI2022103506} proved the following lemma establishing the degree conditions for the existence of a $\mathcal{B}S_\ell$.

% \begin{lem}[Khormali and Palmer~\cite{KHORMALI2022103506}]\label{lem3.1}
% \ \ \ \ \

% \textbf{(i)} Fix integers $\ell>r\geq2$ and let $\cH$ be an $r$-graph. If $x$ is a vertex of degree $d(x)>\binom{\ell-1}{r-1}$ in $\cH$, then $\cH$ contains a $\mathcal{B}S_\ell$ with center $x$.

% \textbf{(ii)} Fix integers $r\geq2$ and $\ell\leq r$ and let $\cH$ be an $r$-graph. If $x$ is a vertex of degree $d(x)>\ell-1$ in $\cH$, then $\cH$ contains a $\mathcal{B}S_\ell$ with center $x$.
% \end{lem}


% Now let us start with the proof of Theorem~\ref{thm: berge matching and berge bigraph} that we restate here for convenience.

% \begin{thm*}
%     Fix integers $s$ and $r\geq2$. Let $G$ be a bipartite graph with $p=p(G)$ and $n$ be sufficiently large.

%     \textbf{(i)} If $r< p\leq s$, then $\ex_r(n,\mathcal{B}M_{s+1}\cup\mathcal{B}G)= \binom{p-1}{r-1}n+O(1)$.

%     %\textbf{(ii)} If $p\leq r\leq s$, then......

%     \textbf{(ii)} If $p\leq s<r-1$, then $\ex_r(n,\mathcal{B}M_{s+1}\cup\mathcal{B}G)\leq\max\big\{s,\binom{2s+1}{r}\big\}$.

%     \textbf{(iii)} If $p>s$, then $\ex_r(n,\mathcal{B}M_{s+1}\cup\mathcal{B}G)= \ex_r(n,\mathcal{B}M_{s+1})$.
% \end{thm*}

% \begin{proof}[\bf Proof] %of Theorem~\ref{thm: berge matching and berge bigraph}.]
%     {\bf (i)} For the lower bound, we consider the following $r$-graph: take a set $A$ of order $p-1$ and a set $B$ of order $n-p+1$, and add each hyperedge that contains exactly $r-1$ vertices of $A$. Clearly, this resulting $r$-graph is $\mathcal{B}M_{p}$-free as $|A|<p$, and thus $\mathcal{B}M_{s+1}\cup\mathcal{B}G$-free as $p\leq s$ and $\mathcal{B}G$ contains a $\mathcal{B}M_{p}$.

%     We now continue with the upper bound. Let $\cH$ be an $n$-vertex $\mathcal{B}M_{s+1}\cup\mathcal{B}G$-free $r$-graph with maximum number of hyperedges.
%     Let $\cM$ be a largest Berge matching in $\cH$, and let $U$ be the core of $\cM$. Then $|U|\le 2s$ as $\cH$ is $\mathcal{B}M_{s+1}$-free.
%     % Then $|U| \le 2s$ as $\cH$ is $\mathcal{B}M_{s+1}$-free.
%     Observe that each hyperedge not belonging to $\cM$ intersects $U$ in at least $r-1$ vertices. Otherwise, a larger Berge matching would appear. Now we delete these at most $s$ hyperedges from $\cM$. Denote by $\cH'$ this resulting $r$-graph. Thus, $\cH'$ contains only two types of hyperedges: those containing exactly $r-1$ vertices from $U$, and those containing exactly $r$ vertices from $U$. The number of hyperedges of the first type is at most $\binom{|U|}{r-1}(n-|U|)$, and the number of hyperedges of the second type is at most $\binom{|U|}{r}$.

%     Let $X=\{v\in V(\cH)\backslash U\,|\,d_{\cH'}(v)>\binom{p-1}{r-1}\}$ and $Y=V(\cH)\backslash (U\cup X)$. By Lemma~\ref{lem3.1} {\bf(i)}, for any $v\in X$, there exists a $\mathcal{B}S_p$ in $\cH'$ with center $v$ such that every vertex other than $v$ belongs to $U$. In particular, the core of this $\mathcal{B}S_p$ is contained in $U$, except for the vertex $v$. Note that there are at most $\binom{|U|}{p}$ $p$-sets in $U$. If $|X|\geq\binom{|U|}{p}(|V(G)|-p)$, then by pigeonhole principle, there exists a $p$-set $U'\subseteq U$ that forms the $p$ leaf vertices in the cores of at least $|V(G)|-p$ Berge copies of $S_p$.
%     This implies that $\cH$ contains a $\cB K_{p,|G|-p}$, and therefore a $\mathcal{B}G$, a contradiction. Thus, $|X|<\binom{|U|}{p}(|V(G)|-p)$. Therefore, the number of hyperedges of $\cH$ is at most
%     \begin{eqnarray*}
%     s+\binom{|U|}{r}+\binom{|U|}{r-1}|X|+\binom{p-1}{r-1}|Y|\leq \binom{p-1}{r-1}n+O(1).
%     \end{eqnarray*}

%     {\bf (ii)} If $r-1>s$, then by Theorem~\ref{thm: Khormali Palmer},
%     \[
%         \ex_r(n,\mathcal{B}M_{s+1}\cup \mathcal{B}G)\leq \ex_r(n,\mathcal{B}M_{s+1})\leq \max\left\{s,\binom{2s+1}{r}\right\}.
%     \]

%     {\bf (iii)} By the definition of $p$, $G$ contains a matching $M_p$. Then any Berge copy of $G$ contains a Berge copy of $M_p$. Since $p \ge s+1$, we obtain $\ex_r(n,\mathcal{B}M_{s+1}\cup\mathcal{B}G)= \ex_r(n,\mathcal{B}M_{s+1})$.
% \end{proof}


% Recall that for $r\geq3$, let $\cG_s^{r-1}$ be an $s$-vertex $(r-1)$-graph with the maximum number of hyperedges such that $\partial\cG_s^{r-1}$ is $K_{p,q}$-free.


% Let us continue with the proof of Theorem~\ref{thm: berge matching and berge Kpq} that we restate here for convenience.

% \begin{thm*}%\label{lem4}
%     Fix integers $s$ and $r\geq 2$. Let $n$ be sufficiently large.

%     \textbf{(i)} If $p\leq r \leq \min\{s,p+q-1\}$, then
%     \[
%         \ex_r(n,\mathcal{B}M_{s+1}\cup\mathcal{B}K_{p,q})=(p-1)n+O(1).
%     \]

%     \textbf{(ii)} If $r=p+q \leq s$, then
%     \[
%     \ex_r(n,\mathcal{B}M_{s+1}\cup\mathcal{B}K_{p,q}) = \min\{p-1,e(\cG_s^{r-1})\}\cdot n+O(1).
%     \]
%     % Moreover, if $(p-1)(p+q-1)\leq s$, then $$\ex_r(n,\mathcal{B}M_{s+1}\cup\mathcal{B}K_{p,q})=(p-1)n+O(1).$$

%     \textbf{(iii)} If $p+q<r \leq s$, then
%     \[
%         \ex_r(n,\mathcal{B}M_{s+1}\cup\mathcal{B}K_{p,q})\leq s+\binom{2s}{r}+\binom{2s}{r-1}(pq-1).
%     \]
% \end{thm*}

% \begin{proof}[\bf Proof]
%     We first prove the lower bounds in \textbf{(i)} and \textbf{(ii)}.
%     For the lower bound in \textbf{(i)}, we consider the following $r$-graph.
%     Let $A$ be a fixed $r$-set with vertices $\{a_1,a_2,\ldots,a_r\}$. Let $A_i = A\backslash \{a_i\}$ for $i=1,2,\ldots,p-1$.
%     For each of the remaining $n-r$ vertices, add all hyperedges that contain $A_i$ and the vertex, $i=1,2,\ldots,p-1$.
%     % From a fixed $r$-set, choose $p-1$ $(r-1)$-sets, and for each of them add hyperedges by joining it with every vertex from the remaining $n-r$ vertices.
%     Clearly, this resulting $r$-graph is $\mathcal{B}M_{s+1}$-free as $r \leq s$. Observe that in a $\mathcal{B}K_{p,q}$, there are at least $p+q$ vertices whose degree is at least $p$. However, in this resulting $r$-graph, the number of vertices of degree at least $p$ is at most $r \le p+q-1$. This implies that this resulting $r$-graph is $\mathcal{B}K_{p,q}$-free.

%     For the lower bound in \textbf{(ii)}, we consider the following $r$-graph.
%     %Recall that $\cG_s^{r-1}$ is the $s$-vertex $(r-1)$-graph with the maximum number of hyperedges such that $\partial\cG_s^{r-1}$ is $K_{p,q}$-free.
%     If $e(\cG_s^{r-1}) \le p-1$, let $\cA = E(\cG_s^{r-1})$. Otherwise, let $\cA$ be a collection consisting of exactly $p-1$ hyperedges of $\cG_s^{r-1}$.
%     For each of the remaining $n-s$ vertices and each $A \in E(\cA)$, add the hyperedge containing the vertex and $A$.
%     Note that $|\cA|=\min\{p-1, e(\cG_s^{r-1})\}$ and the resulting $r$-graph has exactly $(n-s)|\cA|$ hyperedges.
%     Clearly, this resulting $r$-graph is $\mathcal{B}M_{s+1}$-free as $\cG_s^{r-1}$ has only $s$ vertices.
%     Observe that only the vertices in $V(\cG_s^{r-1})$ have degree greater than $p$.
%     Therefore, the core of any $\mathcal{B}K_{p,q}$ is contained only in $\cG_s^{r-1}$. However, by the definition of $\cG_s^{r-1}$, $\partial\cG_s^{r-1}$ is $K_{p,q}$-free. This implies that this resulting $r$-graph is $\mathcal{B}K_{p,q}$-free.
%     % In particular, if $(p-1)(p+q-1)\leq s$, then there exists an $(r-1)$-uniform matching $M_{p-1}^{r-1}$. Since $\partial M_{p-1}^{r-1}$ is $K_{p,q}$-free, $e(\cG_s^{r-1})\geq p-1$.

%     We now prove the upper bounds in \textbf{(i)}, \textbf{(ii)} and \textbf{(iii)} together. Let $\cH$ be an $n$-vertex $\mathcal{B}M_{s+1}\cup\mathcal{B}K_{p,q}$-free $r$-graph.
%     Similar to the proof of Theorem~\ref{thm: berge matching and berge bigraph}, let $\cM$ be a largest Berge matching in $\cH$, and let $U$ be the core of $\cM$.
%     Now we delete at most $s$ hyperedges from $\cM$ and denote by $\cH'$ this resulting $r$-graph. Then $\cH'$ contains only two types of hyperedges: those containing exactly $r-1$ vertices from $U$, and those containing exactly $r$ vertices from $U$.
%     %The number of hyperedges of the first type is at most $\binom{|U|}{r-1}(n-|U|)$, and the number of hyperedges of the second type is at most $\binom{|U|}{r}$.

%     \smallskip
%     {\bf{Case 1.}} $p\leq r \leq p+q$.
%     \smallskip
%     \smallskip

%     Let $X=\{v\in V(\cH)\backslash U\,|\,d_{\cH'}(v)>p-1\}$ and $Y=V(\cH)\backslash (U\cup X)$. By Lemma~\ref{lem3.1} {\bf(ii)}, for any $v\in X$, there exists a $\mathcal{B}S_p$ in $\cH'$ with center $v$ such that every vertex other than $v$ belongs to $U$. In particular, the core of this $\mathcal{B}S_p$ is contained in $U$, except for the vertex $v$. Note that there are at most $\binom{|U|}{p}$ $p$-sets in $U$. If $|X|\geq \binom{|U|}{p}q$, then by pigeonhole principle, there exists a $p$-set $U'\subseteq U$ that forms the $p$ leaf vertices in the cores of at least $q$ Berge copies of $S_p$.
%     % Then any such $q$ $\mathcal{B}S_p$ form a $\mathcal{B}K_{p,q}$ whose core consists of $U'$ together with the $q$ vertices of these $q$ $\mathcal{B}S_p$ in $X$.
%     This implies that $\cH$ contains a $\mathcal{B}K_{p,q}$, a contradiction. It follows that $|X|<\binom{|U|}{p}q$. Therefore, the number of hyperedges of $\cH$ is at most
%     \begin{eqnarray*}
%     s+\binom{|U|}{r}+\binom{|U|}{r-1}|X|+(p-1)|Y|\leq (p-1)n+O(1).
%     \end{eqnarray*}

%     \smallskip
%     {\bf{Case 2.}} $r = p+q \le s$ and $e(\cG_s^{r-1}) < p-1$.
%     \smallskip
%     \smallskip

%     Let $d:=d(s,r,p,q)$ be a large enough fixed constant.
%     For every $(r-1)$-set $S$ in $U$, set $E_{\cH'}(S)=\big\{e\in E(\cH')\,\big|\, e\cap U=S\big\}$ and $X=\Big\{S\subseteq U\,\big|\,|E_{\cH'}(S)|>d\Big\}$.
%     We claim that $|\bigcup_{S\in X}S|\leq s$.
%     Otherwise, we may greedily construct a $\mathcal{B}M_{s+1}$ such that the maximum independent set of its core is contained in $\bigcup_{S\in X}S$, a contradiction.
%     For convenience, let $|\bigcup_{S\in X}S|:=x\leq s$.
%     Now let $\tilde{\cG}$ denote the $(r-1)$-graph with vertex set $\bigcup_{S\in X}S$ and hyperedge set $X$. We claim that $\partial\tilde{\cG}$ is $K_{p,q}$-free.
%     Otherwise, there exists a copy of $K_{p,q}$ in $\partial\tilde{\cG}$, denoted by $K$. Since each edge of $K$ is contained in some $(r-1)$-set from $X$,
%     by the definition of $X$, we may greedily construct a $\mathcal{B}K_{p,q}$ with $K$ as its core, a contradiction.
%     Thus, we have $|X|=e(\tilde{\cG})\leq e(\cG_{x}^{r-1})\leq e(\cG_s^{r-1})<p-1.$
%     Note that $|E_{\cH'}(S)|\leq n-|U|$ for any $S\subseteq U$.
%     Therefore, the number of hyperedges of $\cH$ is at most
%     \begin{eqnarray*}
%     s+\binom{|U|}{r}+e(\cG_s^{r-1})\cdot(n-|U|)+\left(\binom{|U|}{r-1}-e(\cG_s^{r-1})\right)d=e(\cG_s^{r-1})\cdot n+O(1).
%     \end{eqnarray*}


%     %Let $X=\{v\in V(\cH)\backslash U \mid d_{\cH'}(v)> e(\cG_s^{r-1}) \}$ and $Y=V(\cH)\backslash (U\cup X)$. Assume $U = \{u_1,v_1,u_2,v_2,\ldots,u_t,v_t\}$ and $\cM = \{e_1,e_2,\ldots,e_t\}$ such that $u_i,v_i\in e_i$ for $1\leq i\leq t$.
%     %Let $Z_i = \{v \in V(\cH)\backslash U\mid \exists e_1,e_2 \in E(\cH)\setminus E(\cM), v, u_i\in e_1, v,v_i\in e_2\}$. By the maximality of $\cM$, $|Z_i| \le 1$ for all $1\leq i\leq t$. Let $Z = \bigcup_{i=1}^{t}Z_i$,  $X' = X\setminus Z$.
%     %For every $v \in X'$, it contains at most one from $\{u_i,v_i\}$ for each $1 \le i \le t$. If $X' = \emptyset$. Then the number of hyperedges of $\cH$ is at most
%     %\[
%     %    s + \binom{|U|}{r} + t\binom{|U|}{r-1} + e(\cG_s^{r-1})|Y|\leq e(\cG_s^{r-1})n+O(1).
%     %\]
%     %Otherwise, for each pair $\{u_i,v_i\}$, we may assume there is no hyperedge containing $u_i$ and $x \in X'$. That is, for every hyperedge containing $x' \in X'$, it must lie in $V' \cup \{x'\}$ where $V' = \{v_1,v_2,\ldots,v_t\}$.

%     %For any vertex $v \in X\setminus Z$, noting that $V(\cL_v)\subseteq V'$, there exists a $K_{p,q}$ in $\partial \cL_v$. Note that there are at most $\binom{|V'|}{p+q}\binom{p+q}{p}$ copies of $K_{p,q}$ in $V'$. If $|X\setminus Z|\geq \binom{|V'|}{p+q}\binom{p+q}{p}pq$, then by pigeonhole principle, there exists a $K_{p,q}$ in $V'$ such that, there exists at least $pq$ vertices $x_1,\ldots,x_{pq}$ satisfying every $\partial \cL_{x_i}$ contains this $K_{p,q}$. Then for each edge in the $K_{p,q}$, we can find a distinct hyperedge in $\cH$ containing this edge and thus, resulting in a $\mathcal{B}K_{p,q}$, a contradiction. It follows that $|X|<\binom{s}{p+q}\binom{p+q}{p}pq + s$. Therefore, the number of hyperedges of $\cH$ is at most
%     %\begin{eqnarray*}
%     %s+\binom{|U|}{r}+\binom{|U|}{r-1}|X|+e(\cG_s^{r-1})|Y|\leq e(\cG_s^{r-1})n+O(1).
%     %\end{eqnarray*}


%     \smallskip
%     {\bf{Case 3.}} $p+q<r \leq s$.
%     \smallskip
%     \smallskip

%     We claim that any $(r-1)$-set $S \subseteq U$ is contained in at most $pq-1$ hyperedges. Indeed, If there exists such an $(r-1)$-set $S \subseteq U$, then we may greedily construct a $\mathcal{B}K_{p,q}$ such that its core is contained in $S$. Thus, the number of hyperedges in $\cH'$ that contain exactly $r-1$ vertices from $U$ is at most
%     $\binom{|U|}{r-1}(pq-1)$. Therefore, the number of hyperedges of $\cH$ is at most
%     \begin{eqnarray*}
%     s+\binom{|U|}{r}+\binom{|U|}{r-1}(pq-1)\leq s+\binom{2s}{r}+\binom{2s}{r-1}(pq-1).
%     \end{eqnarray*}
%     This completes the proof.
% \end{proof}



% We now begin with the proof of Theorem~Theorem~\ref{prop2.8} that we restate here for convenience.

% \begin{prop*}
%     Let integers $r\geq2$ and $n$ be sufficiently large.

%     \textbf{(i)} If $p=r=s+1$, then $\ex_r(n,\mathcal{B}M_{s+1}\cup\mathcal{B}K_{p,q})=\ex_r(n,\mathcal{B}M_{s+1})$.

%     \textbf{(ii)} If $p<r=s+1\leq p+q$, then $\ex_r(n,\mathcal{B}M_{s+1}\cup\mathcal{B}K_{p,q})= n+O(1)$.

%     \textbf{(iii)} If $p+q<r=s+1$, then $\ex_r(n,\mathcal{B}M_{s+1}\cup\mathcal{B}K_{p,q})\leq s+\binom{2s}{r}+\binom{2s}{r-1}(pq-1)$.
% \end{prop*}

% Here we make a slight adjustment to the proof of the upper bound in Theorem~\ref{thm: berge matching and berge Kpq}.

% \begin{proof}[\bf Proof]
%     %It is easy to see that $M_{s+1}\subseteq K_{p,q}$ as $p=s+1$, and therefore \textbf{(i)} follows.
%     Obviously, \textbf{(i)} follows directly from Theorem \ref{thm: berge matching and berge bigraph} \textbf{(iii)}.
%     For the lower bound in \textbf{(ii)}, we consider the following $r$-graph.
%     For a fixed $(r-1)$-set, add hyperedges by joining it with every vertex from the remaining $n-r+1$ vertices.
%     Clearly, this resulting $r$-graph is $\mathcal{B}M_{s+1}$-free as $r-1<s+1$.
%     Observe that in a $\mathcal{B}K_{p,q}$, there are at least $p+q$ vertices whose degree is at least $p$.
%     However, in this resulting $r$-graph, the number of vertices of degree at least $p$ is at most $r-1$.
%     This implies that this resulting $r$-graph is $\mathcal{B}K_{p,q}$-free as $r-1<p+q$.

%     We now prove the upper bounds in \textbf{(ii)} and \textbf{(iii)} together.
%     Let $\cH$ be an $n$-vertex $\mathcal{B}M_{s+1}\cup\mathcal{B}K_{p,q}$-free $r$-graph.
%     Similar to the proof of Theorem~\ref{thm: berge matching and berge bigraph}, let $\cM$ be a largest Berge matching in $\cH$, and let $U$ be the core of $\cM$.
%     Now we delete at most $s$ hyperedges from $\cM$ and denote by $\cH'$ this resulting $r$-graph. Then $\cH'$ contains only two types of hyperedges: those containing exactly $r-1$ vertices from $U$, and those containing exactly $r$ vertices from $U$.
%     %The number of hyperedges of the first type is at most $\binom{|U|}{r-1}(n-|U|)$, and the number of hyperedges of the second type is at most $\binom{|U|}{r}=O(1)$.

%     Suppose $p<r=s+1\leq p+q$.
%     Let $d:=d(s,r,p,q)$ be a large enough fixed constant.
%     For every $(r-1)$-set $S$ in $U$, set $E_{\cH'}(S)=\big\{e\in E(\cH')\,\big|\, e\cap U=S\big\}$ and $X=\Big\{S\subseteq U\,\big|\,|E_{\cH'}(S)|>d\Big\}$.
%     We claim that $|X|\leq 1$.
%     Otherwise, there are two $(r-1)$-sets $S_1$ and $S_2$ such that $|E_{\cH'}(S_1)|>d$ and $|E_{\cH'}(S_2)|>d$.
%     Since $|S_1\cup S_2|\geq r=s+1$, we may greedily construct a $\mathcal{B}M_{s+1}$ such that the maximum independent set of its core is contained in $S_1\cup S_2$, a contradiction.
%     Note that $|E_{\cH'}(S)|\leq n-|U|$ for any $S\subseteq U$.
%     Therefore, the number of hyperedges of $\cH$ is at most
%     \begin{eqnarray*}
%     s+\binom{|U|}{r}+(n-|U|)+\left(\binom{|U|}{r-1}-1\right)d=n+O(1).
%     \end{eqnarray*}

%     If $p+q<r=s+1$, an argument analogous to that of Theorem~\ref{thm: berge matching and berge Kpq} \textbf{(iii)} yields $e(\cH)\leq s+\binom{2s}{r}+\binom{2s}{r-1}(pq-1)$.
%     This completes the proof.
% \end{proof}



% \section{\normalsize Proofs of Lemma~\ref{lem3} and Theorem~\ref{lem4}}\label{sec: lems}
% %In this section, we use a slightly different notation. We will use $\cF$ to denote a finite family of graphs.

% We first introduce Gallai-Edmonds Structure Theorem on bipartite graphs.
% For a graph $G$, we denote by $D$ the set of all vertices in $G$ that are not covered by at least one maximum matching of $G$.
% Let $A:=N_G(D)$ and $C= V(G)\setminus (D\cup A)$.
% %Let $A$ be the neighborhood of $D$, i.e., the set of vertices in $V(G)\setminus D$ adjacent to at least one vertex in $D$. Finally, let $C= V(G)\setminus (D\cup A)$.
% The partition $(D,A,C)$ is called the \emph{Gallai-Edmonds  decomposition} or \emph{G-E decomposition} shortly.
% Here we use the following result (see Theorem 3.2.4 in~\cite{lovasz1986plummer}) to prove Lemma \ref{lem3}, which extends the result in~\cite{GERBNER2020103082}.

% \begin{thm}[Lov\'{a}sz and Plummer \cite{lovasz1986plummer}]\label{Gallai}
% For a bipartite graph $X = (X_1, X_2) $ and $i = 1, 2$, let $A_i = A \cap X_i $, $C_i = C \cap X_i $ and $D_i = D \cap X_i $, where $ A,  C $, and $ D$ are  the G-E decomposition for $X $. Then

% \textbf{(i)} $ D_1\cup D_2 $ is an independent set of $X$.

% \textbf{(ii)} The subgraph $X[C_1\cup C_2] $ has a perfect matching.

% \textbf{(iii)} $N_X(D_1) = A_2$ and $ N_X(D_2) = A_1$.

% \textbf{(iv)} For any maximum matching $M$ of $X$, it contains a perfect matching of  $ X[C_1 \cup C_2]$, and contains all vertices of $A_1$~(resp. $A_2$) with neighbors in $D_2$~(resp. $D_1$).
% %\item [\rm{(v)}] the subgraphs induced by $A_1 \cup D_2$ and $A_2 \cup D_1$ have positive surplus when viewed from $A_1$ and $A_2$ respectively.
% \end{thm}

% Let $\cH$ be an $r$-graph. Define an \textit{auxiliary bipartite graph} $X$ with parts $X_1$ and $X_2$, where $X_1$ is the hyperedges of $\cH$ and $X_2$ is the collection of all pairs of vertices contained in some hyperedges of $\cH$.
% A vertex $x\in X_1$ is joined to $y\in X_2$ if and only if the hyperedge $x$ contains the pair $y$.


% Now let us start with the proof of Lemma~\ref{lem3} that we restate here for convenience.

% \begin{lem*}
% Let $\mathcal{F}$ be a finite family of graphs. Then
% \begin{eqnarray*}
% {\rm{ex}}_r(n,\mathcal{B}\mathcal{F})\leq \max\Big\{\varepsilon(r,G_{{\rm red{\text -}blue}})\,\big|\, G\ {\rm is}\ {\rm an}\ n{\rm {\text -}vertex}\ \mathcal{F}{\rm {\text -}free}\ {\rm graph}\Big\}.
% \end{eqnarray*}
% \end{lem*}

%  % (see Lemma 2 in~\cite{GERBNER2020103082}).

% \begin{proof}[{\bf{Proof}}]
% Let $\cH$ be a Berge-$\mathcal{F}$-free $r$-graph of order $n$, and let $X=(X_1,X_2)$ be an auxiliary bipartite graph defined as above.
% For $i = 1, 2$, we let $A_i = A \cap X_i $, $C_i = C \cap X_i $ and $D_i = D \cap X_i $, where $ A,  C $, and $ D$ are the G-E decomposition of $X$.
% Let $M$ be a maximum matching in $X$.
% By Theorem~\ref{Gallai} \textbf{(iv)}, there exist vertices $D'_1\subseteq  D_1$ and $D'_2\subseteq D_2$ respectively that are uncovered by $M$.
% The vertices in $X_2\setminus D'_2 $ correspond to a collection of pairs on  $V(\cH)$.
% Let $G^*$ be the graph on $V(\cH)$ with edges $X_2\setminus D'_2$.
% Clearly, $e(G^*) = |X_2\setminus D'_2| $ and $|V(G^*)| \leq n$.

% We assert that $G^*$ is $\mathcal{F}$-free. If  \(G^*\) contains some \(F \in \mathcal{F}\),
% then by the definition of $X$, we see that for each pair $\{v_1v_2\} \in E(F) \subseteq X_2\setminus D'_2 $, there exists a distinct hyperedge \(e \in X_1\setminus D'_1 \) in $\cH$ such that $\{e, \{v_1,v_2\}\} \in M$, which implies that $\cH$ contains a Berge copy of $F$, a contradiction. Thus the assertion holds.


% We next find a red-blue coloring of $G^*$ that satisfies the statement of
% the lemma. If $D'_1=\emptyset$, then $|D_1|=|A_2|$ by Theorem \ref{Gallai} \textbf{(iii)} and \textbf{(iv)}.  We
% color all edges of $G^*$ red. By Theorem \ref{Gallai} \textbf{(ii)} and \textbf{(iv)}, one can see that
% \begin{align*}
% e(\cH) = |X_1| = |C_1|+|A_1|+|D_1| = |C_2|+|D_2\setminus D'_2|+|A_2|=|X_2\setminus D'_2| = e(G^*).
% \end{align*}
% So  we may assume that $ D'_1\neq \emptyset$. Let \(G^*_{\text{red-blue}}\) be the red-blue graph obtained by coloring the edges of \(G^*\) corresponding to vertices in $C_2\cup (D_2\setminus D'_2)$ red, and those corresponding to vertices in $A_2$ blue.   By Theorem \ref{Gallai} \textbf{(ii)}, \textbf{(iii)}, and \textbf{(iv)}, we have $|C_1\cup A_1|=|C_2\cup (D_2\setminus D'_2)|=e(G^*_{\text{red}})$,  and $N_{X}(D_1)=A_2$.  Note that each vertex in $D_1$ is adjacent to exactly \(\binom{r}{2}\) vertices in $A_2$. So each vertex in $D_1$ corresponds to a blue \(K_r\) in \(G^*_{\text{red-blue}}\), which implies that  $|D_1|\leq \mathcal{N}(K_r, G^*_{\text{blue}})$. Therefore,
% \begin{align*}
% e(\cH) = |X_1| = |C_1|+|A_1| + |D_1| \leq e(G^*_{\text{red}}) + \mathcal{N}(K_r, G^*_{\text{blue}}).
% \end{align*}
% \end{proof}


% We now continue with the proof of Theorem~\ref{lem4} that we restate here for convenience.

% \begin{thm*}
% Let $\mathcal{F}$ be a finite family of graphs and $f := f(n)$ such that $ {\rm ex}(n',K_{r-1},\mathcal{F}^-) \le f n$ for every $n$, then
% \begin{eqnarray*}
% {\rm{ex}}_r(n,\mathcal{B}\mathcal{F})\leq \max\Big\{\frac{2f}{r},1\Big\}\cdot{\rm{ex}}(n,\mathcal{F}).
% \end{eqnarray*}
% \end{thm*}

% We will follow the strategy of the proof in~\cite{GERBNER2020103082} by Gerbner, Methuku and Palmer.

% \begin{proof}[{\bf Proof}]
%     Let $G_{{\rm red{\text -}blue}}$ be an $n$-vertex $\mathcal{F}$-free red-blue graph.
% By Lemma~\ref{lem3}, it is sufficient to prove that
% \[
%     \varepsilon(r,G_{{\rm red{\text -}blue}})\leq \max\Big\{\frac{2f}{r},1\Big\}\cdot{\rm{ex}}(n,\mathcal{F}).
% \]

% Write $e(G_{{\rm blue}})=m$, then
% \begin{align}\label{ieq1}
% e(G_{{\rm red}})=e(G_{{\rm red{\text -}blue}})-m\leq {\rm ex}(n,\mathcal{F})-m.
% \end{align}
% Since the underlying graph $G$ is $\mathcal{F}$-free, the neighborhood of every vertex $v\in V(G_{{\rm blue}})$ in $G_{{\rm blue}}$ is $\mathcal{F}^-$-free. Note that an $\mathcal{F}^-$-free graph on $|N_{G_{{\rm blue}}}(v)|$ vertices contains at most
% \begin{align*}
% {\rm ex}(|N_{G_{{\rm blue}}}(v)|,K_{r-1},\mathcal{F}^-)\leq f|N_{G_{{\rm blue}}}(v)|=fd_{G_{{\rm blue}}}(v)
% \end{align*}
%  copies of $K_{r-1}$. This implies that $v$ is contained in at most $fd_{G_{{\rm blue}}}(v)$ copies of $K_r$ in $G_{\rm{blue}}$.
% Summing the number of copies of $K_r$ containing each vertex in $V(G_{{\rm blue}})$, we count each $K_r$ exactly $r$ times. So by $\sum_{v\in V(G_{blue})}d(v)=2m$, one can see that
% \begin{align}\label{ieq2}
% \mathcal{N}(K_r,G_{{\rm blue}})&\leq \frac{1}{r}\sum_{v\in V(G_{{\rm blue}})}{\rm ex}(|N_{G_{{\rm blue}}}(v)|,K_{r-1},\mathcal{F}^-) \notag\\
% &\leq \frac{1}{r}\sum_{v\in V(G_{{\rm blue}})}fd_{G_{{\rm blue}}}(v) \notag\\
% &= \frac{2fm}{r}.
% \end{align}
% Thus, by (\ref{ieq1}) and (\ref{ieq2}), we have
% \begin{eqnarray*}
% \varepsilon(r,G_{{\rm red{\text -}blue}})&\leq& {\rm ex}(n,\mathcal{F})-m+\frac{2fm}{r} \\
% &\leq& \max\Big\{\frac{2f}{r},1\Big\}\cdot\big({\rm ex}(n,\mathcal{F})-m\big)+\max\Big\{\frac{2f}{r},1\Big\}\cdot m \\
% &=& \max\Big\{\frac{2f}{r},1\Big\}\cdot {\rm ex}(n,\mathcal{F}).
% \end{eqnarray*}
% This completes the proof.
% \end{proof}



\bigskip

\textbf{Funding}:

The research of Wang is supported by the China Scholarship Council (No.~202506210200) and the National Natural Science Foundation of China (Grant 12571372).

The research of Yang is supported by the National Natural Science Foundation of China (Nos.~12401464 and 12471334).

The research of Zhao is supported by the China Scholarship Council (No.~202506210250) and the National Natural Science Foundation of China (Grant 12571372).

The research of Bai is supported by the National Natural Science Foundation of China (Nos. 12131013 and 12471334), the Shaanxi Fundamental Science Research Project for Mathematics and Physics (No. 22JSZ009) and the China Scholarship Council (No. 202406290002).

The research of Zhou is supported by the China Scholarship Council (No.~202406890088) and the National Natural Science Foundation of China (Nos.~12271337 and 12371347).


\bigskip

\noindent{\bf{Declaration of interest}}

The authors declare no known conflicts of interest.

% ---------------- bib -------------------------------
\bibliography{ref.bib}
\bibliographystyle{wyc4}



\end{document}

